\chapter{Metode Penelitian}

%% VERSI BAHASA INDONESIA. Konvensi peristilahan dijelaskan pada kepala chapter-1.tex.
%% Kode TikZ, label, sitasi, dan persamaan dipertahankan identik dengan versi Inggris;
%% hanya prosa, caption, dan kepala tabel yang diterjemahkan.

\section{Gambaran Umum Penelitian}

Penelitian ini mengikuti \emph{design science research} (DSR), yaitu paradigma untuk
penelitian yang kontribusinya berupa artefak yang dibangun, bukan hipotesis yang diuji. Jika
penelitian keperilakuan berupaya menjelaskan atau memprediksi, \emph{design science} berupaya
memperluas kapabilitas dengan membangun dan mengevaluasi sebuah artefak dan pengetahuan
tentang masalah diperoleh \emph{melalui} proses pembangunan tersebut \cite{hevner2004}. Model
proses spesifik yang diikuti adalah \emph{design science research methodology} dari Peffers
\textit{et al.} \cite{peffers2007}, yang terdiri atas enam aktivitas: identifikasi masalah dan
motivasi, penetapan tujuan solusi, perancangan dan pengembangan, demonstrasi, evaluasi, serta
komunikasi.

Artefak dalam penelitian ini adalah sistem rekomendasi klinis (CRS) berbasis
\emph{knowledge graph} yang FHIR-native untuk penyakit jantung koroner kronis (CCD);
evaluasinya berupa penilaian keselarasan, kesepakatan, keamanan, dan \emph{faithfulness}
pada Subbab~\ref{sec:eval-method}. Keenam aktivitas DSR dipetakan pada skripsi ini sebagai berikut. Identifikasi masalah dan
tujuan berada pada Bab~I. Perancangan dan pengembangan berupa empat tahap di bawah.
Demonstrasi berupa \emph{grounded briefing} yang dihasilkan untuk masing-masing pasien.
Evaluasi berada pada Bab~IV dan komunikasinya adalah skripsi ini sendiri. Penelitian disusun dalam
empat tahap:

\begin{enumerate}
    \item \textbf{Pembangkitan data.} Rekam medis pasien sintetis dihasilkan dengan Synthea
    dan diekspor sebagai \emph{bundle} HL7 FHIR R4.
    \item \textbf{Pembangunan \emph{knowledge graph}.} \emph{Bundle} FHIR dimuat ke dalam
    \emph{knowledge graph} Neo4j, yang kemudian disatukan dengan panduan AHA/ACC CCD 2023 dan
    korpus literatur biomedis yang diambil dari referensi panduan klinis tersebut melalui lapisan konsep
    berbasis UMLS.
    \item \textbf{Implementasi sistem rekomendasi.} Lapisan kueri mencocokkan setiap pasien
    dengan rekomendasi panduan klinis yang berlaku, menerapkan \emph{observation gate} pada
    rekomendasi bergantung-observasi berdasarkan nilai terukur, mendeteksi tanda keamanan
    kontraindikasi, memeringkat pilihan yang didukung literatur, dan menghasilkan
    \emph{grounded briefing} dalam bahasa alami.
    \item \textbf{Evaluasi.} Sistem dievaluasi untuk kualitas pembangunan graf, keselarasan
    panduan klinis (dengan ablasi \emph{observation gate}), kesepakatan dengan standar emas berlabel
    manual, deteksi keamanan, serta \emph{faithfulness} ringkasan.
\end{enumerate}

Gambar~\ref{fig:research-flow} merangkum alur tersebut beserta artefak utama yang dihasilkan
setiap tahap. Setiap tahap diuraikan pada subbabnya sendiri: tahap~1 pada
Subbab~\ref{sec:data-source}, tahap~2 pada Subbab~\ref{sec:kg-construction}, tahap~3 pada
Subbab~\ref{sec:dss-method}, dan tahap~4 pada Subbab~\ref{sec:eval-method}.

\begin{figure}[htbp]
\centering
\begin{tikzpicture}[
    node distance = 7mm and 12mm,
    stage/.style = {rectangle, rounded corners=2pt, draw, thick,
                    text width=62mm, align=center, minimum height=11mm,
                    font=\small},
    art/.style   = {rectangle, draw, dashed, text width=32mm, align=center,
                    minimum height=8mm, font=\scriptsize, fill=black!5},
    ar/.style    = {-latex, thick}
]
\node[stage] (s1)
    {\textbf{1. Pembangkitan data}\\[1pt]
     \scriptsize Synthea $\rightarrow$ \emph{bundle} HL7 FHIR R4};
\node[stage, below=of s1] (s2)
    {\textbf{2. Pembangunan \emph{knowledge graph}}\\[1pt]
     \scriptsize FHIR $\rightarrow$ Neo4j; jembatan konsep UMLS;
     lapisan panduan klinis dan literatur};
\node[stage, below=of s2] (s3)
    {\textbf{3. Implementasi sistem rekomendasi}\\[1pt]
     \scriptsize himpunan rekomendasi berlaku, \emph{observation gate}, tanda keamanan,
     pemeringkatan literatur, \emph{grounded briefing}};
\node[stage, below=of s3] (s4)
    {\textbf{4. Evaluasi}\\[1pt]
     \scriptsize struktur graf, keselarasan panduan klinis dan ablasi,
     standar emas, keamanan, \emph{faithfulness}};

\node[art, right=of s1] (a1) {\emph{Bundle} FHIR \\ dan ekspor CSV};
\node[art, right=of s2] (a2) {\emph{Knowledge graph} terpadu};
\node[art, right=of s3] (a3) {Rekomendasi, tanda keamanan, ringkasan};
\node[art, right=of s4] (a4) {Metrik evaluasi};

\draw[ar] (s1) -- (s2);
\draw[ar] (s2) -- (s3);
\draw[ar] (s3) -- (s4);
\draw[ar, dashed] (s1) -- (a1);
\draw[ar, dashed] (s2) -- (a2);
\draw[ar, dashed] (s3) -- (a3);
\draw[ar, dashed] (s4) -- (a4);
\end{tikzpicture}
\caption{Alur penelitian: empat tahap penelitian beserta artefak utama yang dihasilkan
masing-masing tahap.}
\label{fig:research-flow}
\end{figure}

\section{Alat dan Bahan}
\label{sec:tools}

Perangkat lunak dan bahan yang digunakan dalam penelitian ini diuraikan di bawah ini. Versi
dan spesifikasi minimum dinyatakan agar peneliti lain dapat mereproduksi pembangunannya pada
perangkat keras yang setara.

\begin{itemize}
    \item \textbf{Synthea}, pembangkit pasien sintetis \cite{synthea}, \emph{build}
    \texttt{aa0772f}, dijalankan di atas Java~26.0.1, menghasilkan \emph{bundle} HL7 FHIR~R4
    \cite{fhir_r4}. Parameter pembangkitan dicatat pada
    Subbab~\ref{sec:cohort}.
    \item \textbf{Neo4j}, basis data graf edisi Enterprise~2026.05.0, sebagai penyimpan
    \emph{knowledge graph} \cite{neo4j}, dikueri dengan Cypher \cite{cypher}. Instansnya
    dijalankan di bawah Neo4j Desktop~2.2.1 dan dikonfigurasi dengan \emph{heap} JVM sebesar
    4~GB awal dan 8~GB maksimum serta \emph{page cache} 4~GB; grafnya disimpan pada basis
    data bernama \texttt{sihdfhir}.
    \item \textbf{Python} untuk \emph{loader}, jembatan terminologi, kueri pendukung
    keputusan, dan skrip evaluasi, menggunakan \emph{driver} resmi Neo4j.
    Dua interpreter digunakan: Python~3.14.4 untuk \emph{loader}, sistem rekomendasi, dan
    skrip evaluasi (\texttt{neo4j}~6.2.0, \texttt{anthropic}~0.112.0,
    \texttt{biopython}~1.87) serta Python~3.11.9 untuk rangkaian
    pemrosesan bahasa alami, yang oleh \texttt{nmslib} dipatok di bawah 3.12
    (\texttt{scispacy}~0.6.2, \texttt{spacy}~3.7.5, \texttt{numpy}~1.26.4,
    \texttt{scikit-learn}~1.8.0, \texttt{pandas}~3.0.3).
    \item \textbf{scispaCy}, model NLP klinis (\texttt{en\_core\_sci\_md}) dengan
    \emph{entity linker} UMLS \cite{umls} dan MeSH \cite{scispacy}, digunakan untuk ekstraksi
    entitas literatur dan penjembatanan terminologi.
    Modelnya adalah \texttt{en\_core\_sci\_md}~0.5.4; \emph{linker}-nya dibangun terhadap
    subhimpunan UMLS 2022AB (kategori semantik 0, 1, 2, dan 9) serta rilis MeSH 2022. Tiga
    ambang kemiripan digunakan, diperketat seiring presisi yang dituntut tiap tahap: 0,80
    untuk penautan entitas literatur, 0,85 untuk penautan entitas klinis, dan 0,75 untuk
    penjembatanan SNOMED~CT ke UMLS karena nama konsep tertaut lebih andal daripada teks
    bebas.
    \item \textbf{\emph{Large language model}} yang diakses melalui API, digunakan semata-mata untuk
    merumuskan fakta rekomendasi hasil penelusuran menjadi sebuah ringkasan.
    Ringkasan yang dilaporkan di sini dibangkitkan dengan \texttt{claude-sonnet-5}.
    Pilihan model hanya memengaruhi perumusannya:
    penelusuran bersifat deterministik dan selesai sebelum model mana pun dipanggil
    sehingga fakta yang boleh dimuat sebuah ringkasan sudah tetap sebelum pembangkitan
    dimulai.
    \item \textbf{Perangkat keras.} Diperlukan \emph{workstation} yang mampu menampung
    \emph{working set} graf di memori; kendala dominan adalah \emph{page cache} Neo4j selama
    pemuatan massal, bukan CPU.
    Penelitian ini dijalankan pada Intel Core i7 generasi ke-13 dengan model 1365U yang berkecepatan
    1,8~GHz dengan RAM 32~GB di bawah Windows~11 Pro. Basis data sendiri
    mencadangkan 12~GB, yaitu \emph{heap} maksimum 8~GB dan \emph{page cache} 4~GB. Karena
    itu, 16~GB merupakan batas bawah praktis untuk mereproduksi pembangunannya; di bawah
    angka tersebut \emph{page cache} akan saling berebut selama pemuatan massal. RAM 32~GB menyisakan
    ruang sehingga memungkinkan pemuatan dan skrip evaluasi dapat berjalan bersamaan.
\end{itemize}

\section{Sumber Data dan Kohort Penelitian}
\label{sec:data-source}

\subsection{Sumber data}
Data pasien seluruhnya bersifat \emph{sintetis}, dihasilkan oleh Synthea \cite{synthea} dan
diekspor sebagai \emph{bundle} HL7 FHIR R4 \cite{fhir_r4}; tidak ada data pasien nyata yang
digunakan sehingga risiko privasi dan etika dapat ditekan seminimal mungkin. Dalam
\emph{bundle} FHIR, kondisi dan prosedur dikodekan dalam SNOMED~CT \cite{snomedct}, obat dalam
RxNorm \cite{rxnorm}, dan observasi (hasil laboratorium serta tanda vital) dalam LOINC
\cite{loinc}; kode ICD-10-CM/ICD-11 tidak terdapat di dalamnya.

Sebuah \texttt{Bundle} FHIR tidak membawa elemen \texttt{fhirVersion}---hanya
\texttt{CapabilityStatement} yang membawanya dan Synthea tidak menerbitkannya---sehingga
rilis yang digunakan dipastikan dari instans hasil ekspor, bukan diasumsikan. Tiga pembeda
struktural yang saling bebas sama-sama menunjuk pada R4. Pertama, konsep berkode merujuk pada
ruang nama \texttt{terminology.hl7.org/CodeSystem/} tanpa satu pun kemunculan ruang nama
pra-R4 seperti \texttt{hl7.org/fhir/v3/} ataupun \texttt{/v2/}. Kedua, butir klaim menggunakan
\texttt{productOrService}, yang merupakan elemen yang diperkenalkan R4 sebagai pengganti
\texttt{Claim.item.service} pada STU3. Ketiga, pesanan obat muncul sebagai
\texttt{MedicationRequest}, bukan \texttt{MedicationOrder} sebagaimana pada DSTU2. Direktori
keluaran pengekspor (\texttt{fhir/}, bukan \texttt{fhir\_stu3/} atau \texttt{fhir\_dstu2/})
juga menunjukkan konsistensi dengan hal ini. \emph{Bundle} tersebut juga membawa penegasan profil US~Core dan
US~Core hanya didefinisikan terhadap R4.
Bukti pada tingkat instans yang sama juga memastikan versi \emph{implementation guide} yang
dipakai. Himpunan profilnya memuat \texttt{us-core-observation-screening-assessment}, yang
diperkenalkan pada US~Core~6.0.0 untuk menggantikan profil Survey terdahulu---dan profil Survey
itu sendiri tidak ada; \texttt{us-core-condition-encounter-diagnosis} juga muncul, sedangkan profil
tunggal \texttt{us-core-condition} pra-5.0.0 tidak. Synthea menyediakan empat versi
US~Core---3.1.1, 4.0.0, 5.0.1, dan 6.1.0 \cite{synthea_fhir_wiki}---dan hanya 6.1.0 yang berada
pada atau di atas 6.0.0. Karena itu, ekspor tersebut merupakan US~Core~6.1.0. Tidak ada
\texttt{synthea.properties} yang digunakan saat pembangkitan data sehingga
penetapan ini bertumpu pada instans hasil ekspor, setara dengan penetapan rilis di atas.

\subsection{Kohort}
\label{sec:cohort}
Kumpulan data dihasilkan oleh dua proses Synthea yang terpisah, keduanya untuk
negara bagian Massachusetts dengan riwayat sepuluh tahun yang diekspor, lalu digabungkan
menjadi satu korpus (Tabel~\ref{tab:generation}). Proses pertama dibatasi pada modul \texttt{stable\_ischemic\_heart\_disease} agar korpus
memuat kepadatan pasien penyakit sasaran yang memadai. Proses kedua menjalankan seluruh
himpunan modul (\texttt{*}) untuk menyediakan populasi umum. Karena itu, seleksi kohort
harus benar-benar membedakan pasien berdasarkan kode kondisinya, bukan sekadar berdasarkan
proses mana yang membangkitkannya. Kedua proses menghasilkan
$2{.}009 + 3{.}481 = 5{.}490$ pasien secara keseluruhan. Setiap proses mencatat
\emph{seed} acaknya sehingga korpus dapat dibangkitkan ulang secara persis.

\begin{table}[h]
\centering
\caption{Parameter pembangkitan Synthea untuk kedua proses yang menyusun kumpulan data pasien.}
\label{tab:generation}
\begin{tabular}{lll}
\hline
\textbf{Parameter} & \textbf{Proses diperkaya penyakit} & \textbf{Proses umum} \\
\hline
Modul                & \texttt{stable\_ischemic\_} & \texttt{*} (semua) \\
                     & \texttt{heart\_disease} & \\
Jumlah pasien        & 2{.}009 & 3{.}481 \\
\emph{Seed}          & 1784186620874 & 1784186859476 \\
Negara bagian        & Massachusetts & Massachusetts \\
Riwayat (tahun)      & 10 & 10 \\
\emph{Build} Synthea & \texttt{aa0772f} & \texttt{aa0772f} \\
\hline
\multicolumn{3}{l}{\textbf{Total: 5{.}490 pasien}} \\
\hline
\end{tabular}
\end{table}

\subsection{Mengapa pemilihan modul tidak menentukan ukuran kohort}
\label{sec:module-yield}
Ada kemungkinan untuk dapat disalahpahami oleh pembaca bahwa 2{.}009 pasien dari hasil proses Synthea yang diperkaya penyakit sebagai 2{.}009
pasien dengan penyakit jantung koroner kronis. Kenyataannya bukan demikian. Modul Synthea
adalah \emph{mesin keadaan probabilistik}, bukan label pengaktifan kondisi penyakit tertentu. Mengaktifkan modul
\texttt{stable\_ischemic\_heart\_disease} berarti setiap pasien pada proses tersebut
\emph{disimulasikan melewati} modul itu, bukan bahwa setiap pasien pasti mencapai keadaan
diagnosisnya. Perjalanan melalui modul ditentukan oleh usia, jenis kelamin, ras, dan akumulasi
faktor risiko seperti status merokok, diabetes, hipertensi, serta kadar lipid, yang semuanya
juga dihasilkan secara stokastik. Gambar~\ref{fig:sihd-module} menyajikan mesin keadaan modul
tersebut secara lengkap; mekanismenya perlu dinyatakan secara tepat karena mekanisme itulah
yang menjelaskan mengapa hasilnya hanya sebagian dan mengapa pasien yang terdiagnosis berusia
setengah baya.

Modul itu sendiri tidak memuat kondisi usia. Keadaan \texttt{Chance\_of\_IHD} pada modul tersebut merupakan
perulangan berupa tundaan satu bulan yang memiliki penentu untuk melakukan transisi menuju keadaan \emph{onset} berupa atribut \texttt{ihd\_risk} dalam bentuk
peluang, dengan nilai bawaan nol, dan yang selain itu berputar kembali ke dirinya sendiri.
Setiap bulan yang disimulasikan merupakan penarikan yang tak terikat dan seorang pasien dapat berputar
sepanjang hidupnya tanpa pernah bertransisi. Kebergantungan pada usia justru masuk melalui
\texttt{ihd\_risk} itu sendiri, yang dihitung Synthea sebagai risiko kardiovaskular sepuluh
tahun yang dikonversi menjadi \emph{hazard} per bulan lalu dikalikan lima. Pada \emph{build}
yang dipakai di sini, risiko itu berasal dari \emph{ASCVD Pooled Cohort Equations}, bukan dari
Framingham---tanda \texttt{useFramingham} bernilai \texttt{false}---dan
\texttt{ASCVD.ascvd10Year} mengembalikan $-1$ alih-alih sebuah peluang bilamana pasien berusia
di bawah 40 atau di atas 79 tahun. Di luar jendela itu tidak ada \emph{onset} yang dapat terjadi. Di
dalamnya, usia masuk secara kontinu melalui $\ln(\mathrm{usia})$ beserta suku-suku
interaksinya sehingga \emph{hazard} bulanan meningkat sepanjang rentang hidup.

Korpus membuktikan hal ini pada kedua sisi umur tersebut. Dari 876 \emph{onset} yang tercatat, yang
termuda berada pada usia 41 dan mediannya pada usia 60, tanpa satu pun
di bawah 40 dan tepat satu di atas 79---jendela kalkulator itu terreproduksi persis di dalam
data. Karena itu, batas bawah usia~40 bukanlah pertimbangan klinis yang dikodekan ke dalam
modul, melainkan artefak dari rentang tervalidasi kalkulator risikonya dan lebih tepat dibaca
sebagai keterbatasan kohort sintetis daripada sebagai pilihan pemodelan. Kedua proses
tidak berbagi satu pun pasien yang sama dan proses yang diperkaya penyakit sama sekali tidak
mengeluarkan observasi status merokok sementara setiap pasien pada proses umum membawanya
sehingga keberadaan LOINC \texttt{72166-2} menandai asal proses seorang pasien tanpa kebermaknaan
ganda. Dengan penanda itu, proses yang diperkaya penyakit menyumbang 249 dari 876
pasien dan proses umum menyumbang 627---yaitu hasil sebesar 12,4\%
($249/2{.}009$) dan 18,0\% ($627/3{.}481$). \texttt{conditions.csv} milik proses
yang diperkaya penyakit membenarkan angka pertama. Berkas itu memuat tepat
249 baris yang membawa \texttt{414545008}.

Jadi, proses yang dipilih \emph{dengan harapan} menghasilkan penyakit tersebut justru menghasilkan kasus yang secara
proporsional \emph{lebih sedikit} daripada proses yang sekadar mengaktifkan seluruh modul.
Pemilihan modul bukan hanya gagal menentukan ukuran kohort; pemilihan itu menggesernya ke arah
yang berlawanan dengan yang dimaksudkan. Penyebabnya tampak pada ekspor yang sama. Di antara
2{.}009 pasiennya, proses yang diperkaya penyakit hanya mencatat kode kondisi
yang berbeda dan 606 baris kondisi secara keseluruhan, tanpa
hipertensi, diabetes, maupun hiperlipidemia, berbanding 285 kode dan
130.859 baris pada proses umum. \emph{Onset} digerakkan oleh skor risiko ASCVD dan skor
itu justru mengintegrasikan komorbiditas-komorbiditas tersebut: mengaktifkan modul penyakitnya
sendirian berarti membatasi model risikonya dari faktor risiko yang menjadi sandarannya.
Prevalensi penyakit pada korpus Synthea dengan demikian merupakan \emph{luaran} simulasi yang
harus diukur setelah pembangkitan, bukan parameter yang dapat ditetapkan di muka.

Akibatnya, terdapat satu konsekuensi yang nantinya akan dibahas di Bab~IV. Sebanyak 249
pasien dari proses yang diperkaya penyakit di dalam kohort ini membawa satu diagnosis penyakit
jantung iskemik, satu temuan pencitraan jantung yang abnormal, dan pada 108 kasus
riwayat \emph{bypass} koroner---dan tidak ada lagi selain itu. Mereka tidak memiliki komorbiditas
yang dapat dicocoki oleh rekomendasi mana pun dan, selain fraksi ejeksi, tidak memiliki
observasi untuk digerbangi. Karena itu, sekitar 28\% kohort yang dievaluasi jauh lebih miskin
secara klinis daripada pasien penyakit jantung koroner kronis mana pun di kenyataan.

\begin{figure}[h]
\centering
\begin{tikzpicture}[
    node distance = 4.5mm and 10mm,
    st/.style    = {rectangle, rounded corners=2pt, draw, thick, text width=52mm,
                    align=center, minimum height=7mm, font=\scriptsize},
    onset/.style = {st, fill=black!8},
    side/.style  = {rectangle, rounded corners=2pt, draw, dashed, text width=26mm,
                    align=center, minimum height=6mm, font=\scriptsize},
    ar/.style    = {-latex, thick},
    lb/.style    = {font=\tiny, inner sep=1.5pt}
]
\node[onset]                  (init)  {\textbf{Initial}};
\node[onset, below=of init]   (chance){\textbf{Chance\_of\_IHD} --- \emph{Delay} 1 bulan};
\node[onset, below=of chance] (start) {\textbf{IHD\_Start} --- \emph{ConditionOnset}\\
                                       SNOMED 414545008 \emph{Ischemic heart disease}};
\node[st, below=of start]     (delay) {\textbf{Delay\_Until\_Symptoms} --- \emph{Delay} 1 bulan};
\node[st, below=of delay]     (symp)  {\textbf{Symptoms\_Start} $\rightarrow$
                                       \textbf{Dyspnea\_Symptom}\\
                                       nyeri dada; sesak saat aktivitas};
\node[st, below=of symp]      (enc)   {\textbf{Diagnosis\_Encounter} --- \emph{Encounter}
                                       (\emph{wellness})};
\node[st, below=of enc]       (ecg)   {\textbf{ECG / Labs} $\rightarrow$
                                       \textbf{Stress Test}};
\node[st, below=of ecg]       (med)   {\textbf{Medical\_Therapy} --- \emph{Delay} 1 pekan};
\node[st, below=of med]       (presc) {\textbf{Prescribe\_CHD\_Meds}\\
                                       kaskade bergerbang tahun: clopidogrel (1997),
                                       \emph{simvastatin} (1991), metoprolol (1978),
                                       nitrogliserin};
\node[st, below=of presc]     (cabg)  {\textbf{CABG} --- \emph{CallSubmodule}
                                       \texttt{heart/cabg\_sequence}};
\node[st, below=of cabg]      (term)  {\textbf{Terminal}};
\node[side, right=of delay]   (arrest){\textbf{Cardiac\_Arrest} --- \emph{Death}\\
                                       SNOMED 410429000};

\draw[ar] (init)   -- (chance);
\draw[ar] (chance) -- node[lb, right] {$p=\texttt{ihd\_risk}$} (start);
\draw[ar] (start)  -- (delay);
\draw[ar] (delay)  -- node[lb, right] {$0{,}25$} (symp);
\draw[ar] (symp)   -- (enc);
\draw[ar] (enc)    -- (ecg);
\draw[ar] (ecg)    -- node[lb, right] {$0{,}99$ pada tiap langkah} (med);
\draw[ar] (med)    -- (presc);
\draw[ar] (presc)  -- node[lb, right] {$0{,}9$ via Symptoms\_Continue} (cabg);
\draw[ar] (cabg)   -- (term);
\draw[ar, dashed] (delay)  -- node[lb, above] {$0{,}0005$} (arrest);
\draw[ar, dashed] (arrest) |- (term);

\draw[ar] (chance.north west) to[out=180, in=180, looseness=2]
          node[lb, left] {selain itu} (chance.south west);
\draw[ar] (delay.north west)  to[out=180, in=180, looseness=2]
          node[lb, left] {$0{,}7495$} (delay.south west);
\draw[ar] (presc.west)        to[out=180, in=180, looseness=1.5]
          node[lb, left] {$0{,}1$} (med.west);
\end{tikzpicture}
\caption{Modul \texttt{stable\_ischemic\_heart\_disease} sebagai mesin keadaan, pada komit
Synthea \texttt{aa0772f}. Hanya keadaan berlatar abu-abu yang menentukan keanggotaan kohort:
pasien masuk ke \texttt{Chance\_of\_IHD} dan menarik sekali setiap bulan simulasi terhadap
\texttt{ihd\_risk} sehingga sebagian besar pasien berputar hingga meninggal tanpa pernah
mencapai \texttt{IHD\_Start}. Segala hal di bawah wilayah berlatar abu-abu merupakan alur
perawatan, bukan diagnosis, dan hanya berjalan bagi pasien yang memang bertransisi. Peluang
yang tertera adalah peluang yang tercantum apa adanya di dalam modul; \texttt{ihd\_risk}
dipasok dari luar oleh \texttt{CardiovascularDiseaseModule} dan tidak terdefinisi di luar usia
40---79. Demi keterbacaan, dua sisi kecil tidak digambarkan: \texttt{ECG / Labs} dan
\texttt{Stress Test} masing-masing juga mengalirkan $0{,}01$ langsung ke \texttt{CABG}, dengan
melewati terapi medis.}
\label{fig:sihd-module}
\end{figure}

\subsection{Gerbang CCD dan penurunan kohort}
\label{sec:gate}
Keanggotaan kohort karenanya ditetapkan dengan mengueri data yang dibangkitkan, bukan dengan
memercayai konfigurasi proses pembangkitan saja. Seorang pasien masuk ke kohort evaluasi jika graf
memuat, bagi pasien tersebut, sebuah \texttt{:Condition} yang kodenya menjembatani ke konsep
UMLS \texttt{C0151744}---yang oleh lapisan terminologi diresolusi dari SNOMED~CT
\texttt{414545008}, \emph{Ischemic heart disease (disorder)}. Secara konkret, gerbang tersebut
adalah penelusuran Cypher
\texttt{(:Patient)$\leftarrow$[:SUBJECT]$-$(:Condition)$-$[:CODE]$\rightarrow$(:Code)
$\leftarrow$[:SAME\_CONCEPT]$-$(:Concept)} dan gerbang inilah klausa pertama yang dievaluasi
pada kueri penentuan kohort: pasien yang tidak melewatinya tidak pernah dicocokkan dengan
rekomendasi apa pun. Konsep gerbang yang sama digunakan oleh standar rujukan yang menjadi dasar penilaian sistem,
yaitu \emph{silver standard} yang dihasilkan secara otomatis serta didefinisikan dan
dijustifikasi pada Subbab~\ref{sec:oracle}. Dengan demikian, kedua sisi evaluasi keselarasan
dihitung pada kohort yang identik.

Tabel~\ref{tab:cohort-funnel} menelusuri penurunan yang dihasilkannya.

\begin{table}[h]
\centering
\caption{Penurunan kohort yang dapat dievaluasi gerbang dari korpus hasil pembangkitan.}
\label{tab:cohort-funnel}
\begin{tabular}{lrl}
\hline
\textbf{Tahap} & \textbf{Pasien} & \textbf{Ditentukan oleh} \\
\hline
Korpus hasil pembangkitan (kedua proses) & 5{.}490 & konfigurasi proses Synthea \\
\quad proses diperkaya penyakit          & 2{.}009 & \texttt{stable\_ischemic\_heart\_disease} \\
\quad proses umum                        & 3{.}481 & seluruh modul (\texttt{*}) \\
Membawa SNOMED 414545008                 &    876  & luaran simulasi \\
\textbf{Kohort yang dapat dievaluasi gerbang} & \textbf{876} & \textbf{gerbang CCD} \\
\hline
\end{tabular}
\end{table}

Angka 876 dengan demikian adalah jumlah pasien simulasi yang benar-benar mengalami diagnosis
penggerbang (SNOMED~CT~\texttt{414545008}, \emph{Ischemic heart disease (disorder)}), yaitu $16{,}0\%$ dari korpus.

\subsection{Apa yang tidak tercakup oleh gerbang}
Panduan AHA/ACC 2023 mendefinisikan penyakit jantung koroner kronis secara \emph{disjungtif}:
riwayat infark miokardial dan riwayat revaskularisasi koroner sama-sama menjadi penentu CCD di
samping diagnosis penyakit jantung iskemik. Gerbang yang diimplementasikan di sini hanya
menggunakan yang terakhir. Sebanyak 67 pasien lain dalam korpus membawa \emph{History of
myocardial infarction} (SNOMED \texttt{399211009}) atau \emph{History of coronary artery
bypass grafting} (SNOMED \texttt{399261000}) tanpa turut membawa \texttt{414545008}; menurut
definisi panduan klinis, mereka juga termasuk sebagai pasien CCD sehingga kohort yang didefinisikan
secara klinis semestinya berjumlah 943.

Kesenjangan ini bukan kegagalan penjembatanan terminologi. Ketiga kode diresolusi dengan benar ke konsep UMLS yang berbeda: \texttt{414545008} ke
\texttt{C0151744}, \texttt{399211009} ke \texttt{C1275835}, dan \texttt{399261000} ke
\texttt{C1275842}. Pembedaan itu benar secara semantik karena SNOMED~CT memodelkan
\emph{situation with explicit context} (``riwayat'') sebagai entitas yang berbeda dari
\emph{disorder} yang bersesuaian.
Lapisan panduan klinis pun tidak keliru: enam rekomendasi memang secara khusus menargetkan konsep riwayat infark miokardial.
Keterbatasannya berada pada gerbang yang dinyatakan sebagai
konsep \emph{tunggal} sementara panduan klinis menyatakan sebuah disjungsi sehingga keenam
rekomendasi tersebut tidak akan pernah terjangkau oleh 67 pasien itu karena gerbang
dievaluasi sebelum pencocokan rekomendasi dimulai.

Gerbang berkonsep tunggal ini dipertahankan dalam penelitian ini agar angka keselarasan,
kesepakatan, dan keamanan yang dilaporkan seluruhnya menggambarkan satu kohort yang
didefinisikan secara konsisten. Kohort yang dapat dievaluasi gerbang karenanya harus dibaca
sebagai subhimpunan konservatif dari populasi CCD klinis---$876$ dari $943$, atau
$92{,}9\%$---dan sisa $7{,}1\%$ ditinjau kembali di antara keterbatasan pada
Subbab~\ref{sec:ethics}.
Ketiga puluh pasien yang dilabeli secara manual tidak diambil secara acak. Setiap pasien
dari 876 tersebut diberi skor menurut sejauh mana rekamnya dapat menghasilkan ringkasan.
Penskorannya memberi dua poin untuk tiap komorbiditas pembeda, dihitung paling banyak
untuk enam dari sembilan komorbiditas dalam \emph{allow-list}, sehingga menyumbang maksimum
12 poin. Penskorannya juga memberi satu poin untuk tiap kelas obat relevan-keamanan,
dihitung paling banyak untuk tiga dari enam kelas, sehingga menyumbang maksimum 3 poin.
Skor tertinggi yang mungkin dicapai karenanya 15 poin. Skor tersebut kemudian dikenai
penalti atas artefak pembangkit data. Penalti terberat dikenakan pada fraksi ejeksi
ventrikel kiri di luar rentang 20---70\%, yang tidak fisiologis dan justru akan mengundang
pembaca mengkritik Synthea alih-alih sistemnya. Tiga puluh pasien dengan skor tertinggi kemudian diambil.
Prosedurnya deterministik: peringkat~1---18 mencapai skor
tertinggi tersebut dan peringkat~19---30 memperoleh 13 poin, dengan batas potong berada di dalam
kesamaan skor tersebut dan diputus menurut urutan nomor identitas pasien. Sampel ini karenanya memiliki
alasan pemilihan yang jelas, yaitu condong pada rekam yang paling kompleks secara klinis, sehingga kesepakatan yang
dilaporkan pada Subbab~\ref{sec:goldlabel} dimaknai sebagai hasil pengukuran atas
kasus-kasus tersebut, bukan atas pasien rata-rata.

\section{Metode Pembangunan \emph{Knowledge Graph}}
\label{sec:kg-construction}

\emph{Knowledge graph} dibangun secara \emph{EHR-first}: lapisan konsep dicetak dari kode EHR
sebelum korpus literatur dimuat. Resolusi EHR ke UMLS hanya bergantung pada deskripsi kode dan
pada rilis UMLS yang dipakai \emph{linker} sehingga himpunan konsep yang dapat dijangkau dari data pasien---dan bersamanya
setiap sisi \texttt{SAME\_CONCEPT}---tetap utuh betapa pun penelusuran literatur diulang atau
diperluas. Hal ini penting karena gerbang CCD pada Subbab~\ref{sec:gate} melakukan penelusuran menggunakan
sisi-sisi tersebut. Seandainya lapisan konsep dicetak dari literatur, korpus yang berbeda akan
memberikan himpunan konsep yang berbeda dan karenanya kohort yang berbeda pula sehingga
setiap metrik yang dilaporkan dapat berubah tergantung dari penelusuran literatur yang dilakukan. Konsep yang
berasal dari literatur kemudian di-\emph{merge} ke kunci simpul \texttt{umlsCui} yang sama
sehingga memperkaya lapisan konsep tersebut tanpa mendefinisikannya ulang. Pembangunannya berlangsung
sebagai berikut.

\begin{enumerate}
    \item \textbf{Pemuatan FHIR ke graf.} Setiap \emph{resource} FHIR menjadi sebuah simpul,
    setiap rujukan menjadi sebuah relasi bertipe, dan setiap kode terminologi diangkat menjadi
    simpul \texttt{:Code} bersama dengan kunci berupa (\emph{system}, \emph{code}) sehingga semua
    kemunculan sebuah kode merujuk pada satu simpul. Observasi bernilai jamak (misalnya
    tekanan darah) direpresentasikan sebagai simpul komponen agar tidak kehilangan data.
    \emph{Resource} penagihan/administratif dilewati dan hanya \emph{allow-list} observasi LOINC
    yang relevan secara klinis yang dimuat.

    \item \textbf{Jembatan terminologi (EHR ke UMLS).} Kode EHR yang telah dideduplikasi
    diekspor dari graf dan dipetakan ke \emph{concept unique identifier} (CUI) UMLS
    \cite{umls} menggunakan \emph{linker} tiga-lintasan (alias persis, pembangkitan kandidat,
    dan disambiguasi berbasis NLP \cite{scispacy}), menghasilkan lapisan konsep yang terhubung
    ke simpul \texttt{:Code} bersama melalui sisi \texttt{SAME\_CONCEPT}.

    \item \textbf{Lapisan panduan klinis.} Panduan AHA/ACC CCD 2023 dikurasi secara manual menjadi
    81 rekomendasi, masing-masing ditautkan ke konsep yang direkomendasikan,
    dikontraindikasikan, atau ditargetkannya. Tipe sisi diturunkan dari \emph{Class of
    Recommendation} (COR): COR~3 (Harm/No-Benefit) menghasilkan sisi \emph{contraindicates},
    sedangkan COR~1/2a/2b menghasilkan sisi \emph{recommends}. Rekomendasi bertarget
    disjungtif ditandai agar lapisan kueri menggabungkan targetnya dengan OR, bukan AND.

    \item \textbf{Lapisan literatur.} Korpus PubMed yang diekstrak dari panduan klinis (uji klinis
    yang disitasi panduan klinis beserta artikel serupanya) diproses dengan NER klinis scispaCy untuk
    mengekstraksi sebutan entitas, yang lalu ditautkan ke konsep UMLS yang sama. Asosiasi
    obat-kondisi diwujudkan antarsimpul \texttt{:Code} sebagai sisi kemunculan bersama yang
    diperingkat menurut \emph{pointwise mutual information} (PMI) \cite{pmi}, dengan ambang
    minimum jumlah artikel.
    Korpusnya terdiri atas 4.801 abstrak berbeda: 274 uji klinis yang disitasi oleh
    panduan klinis itu sendiri, bersama 4.793 artikel serupa yang ditarik darinya. Kedua himpunan itu
    beririsan cukup lebar---266 abstrak benihnya muncul kembali di antara artikel serupanya
    sendiri dan tiap abstrak hanya disimpan satu kali---sehingga gabungannya lebih kecil
    daripada jumlahnya. Sebuah pasangan obat--kondisi harus disebut bersama dalam
    sekurang-kurangnya tiga abstrak berbeda sebelum layak diperingkat.
\end{enumerate}

\subsection{Model data graf}
\emph{Loader} memetakan FHIR ke graf properti sebagai berikut. Setiap \emph{resource} menjadi
simpul berlabel \texttt{:Resource:<ResourceType>}, dengan atribut skalarnya menjadi bagian dari simpul
tersebut. Elemen yang memuat sebuah rujukan diubah menjadi relasi bertipe yang namanya diturunkan dari
lintasan elemen FHIR; misalnya, \texttt{Observation.subject} menghasilkan sisi
\texttt{SUBJECT}. \emph{Loader} meresolusi ketiga bentuk rujukan yang diterbitkan
Synthea, yaitu pengenal \texttt{urn:uuid:}, rujukan bersyarat
(\texttt{Type?identifier=system\textbar value}), dan rujukan terkandung
(\texttt{\#fragment}).

Label \texttt{:Resource} lebih merupakan urusan teknis daripada urusan makna. Label itu hanya
menyatakan bahwa sebuah simpul dapat dicari melalui \texttt{uid}-nya dan tidak mengeklaim apa
pun tentang apakah simpul itu merupakan \emph{Resource} dalam pengertian FHIR. Setiap bentuk rujukan menyusut menjadi pengenal telanjang (\emph{bare identifier}):
\texttt{Type/id} membuang awalannya dan \texttt{urn:uuid:} memang tidak pernah membawanya.
Akibatnya, ketika sebuah sisi hendak dituliskan, \emph{loader} sudah memegang pengenal
sasarannya, tetapi belum mengetahui tipenya. Karena itu, \emph{loader}
membutuhkan satu label yang dimiliki setiap kemungkinan sasaran dan label inilah yang
berperan demikian. Label yang sama memegang satu-satunya kendala keunikan \texttt{uid}, yang
didukung basis data dengan indeks, sehingga menelusuri sebuah rujukan hanya memerlukan
satu pencarian indeks, bukan pemindaian.

Akibatnya, terdapat dua konsekuensi. Simpul komponen membawa label
\texttt{:Resource} meskipun \texttt{Observation.component} pada FHIR merupakan
\emph{BackboneElement}, bukan \emph{Resource}, karena sisi harus dapat meresolusi ke simpul
tersebut melalui pengenal. Sebaliknya, simpul \texttt{:Code} merupakan pengecualian yang
disengaja: simpul itu memiliki kunci berupa (\texttt{system}, \texttt{code}), bukan berupa \texttt{uid},
tidak pernah menjadi sasaran rujukan \emph{resource}, dan karenanya tidak membawa label
\texttt{:Resource}. Setiap \emph{coding} terminologi diangkat menjadi simpul bersama
\texttt{:Code\{system, code, display\}}, dengan kunci berupa (\texttt{system}, \texttt{code}) di
bawah \emph{constraint} unik sehingga setiap kemunculan sebuah kode---lintas kondisi, prosedur,
obat, dan observasi---merujuk pada satu simpul; simpul \texttt{:Code} bersama inilah jangkar
tempat lapisan konsep UMLS melekat. Sebuah \emph{Coding} dibedakan dari sebuah
\emph{Quantity} (keduanya membawa \texttt{system} dan \texttt{code}) melalui keberadaan
\texttt{value}, yang mencegah satuan ukur terangkat menjadi kode klinis. Observasi bernilai
jamak---terutama panel tekanan darah---dimekarkan menjadi simpul \texttt{:Component} agar
nilai sistolik dan diastolik sama-sama terpelihara. Seluruh penulisan menggunakan
\texttt{MERGE}, yang menjadikan pemuatan bersifat idempoten dan dapat diulang;
\emph{resource} penagihan dan provenans dilewati dan hanya observasi LOINC dalam
\emph{allow-list} yang dimuat.

\subsubsection*{Mengapa beberapa terminologi hadir bersama}
\label{sec:multi-vocab}
Graf yang dimuat memuat 26 \emph{code system} yang berbeda. Hal ini merupakan
sifat FHIR itu sendiri, bukan keanehan generator: spesifikasi dari FHIR mengikat setiap elemen berkode
pada sebuah \emph{value set} dan ikatan tersebut menyebut kosakata tertentu. Synthea
menyatakan kesesuaian itu secara eksplisit---observasi status merokok membawa rujukan
\texttt{meta.profile} ke \texttt{us-core-smokingstatus}---sehingga kosakata yang dipakai
merupakan kosakata yang diwajibkan profil US~Core, bukan pilihan bebas generator. Setiap
kosakata menjawab satu pertanyaan dan tidak berusaha menjawab pertanyaan lainnya. SNOMED~CT
\cite{snomedct} membawa makna klinis---kondisi, prosedur, tipe \emph{encounter} dan
\emph{device}---sebagai ontologi \emph{description logic} dengan \emph{subsumption} yang dapat
dikomputasi. RxNorm \cite{rxnorm} membawa obat, dengan menormalkan penamaan obat lintas
kosakata \emph{proprietary} pada tingkat abstraksi yang eksplisit. LOINC \cite{loinc} membawa
observasi, mengidentifikasi \emph{apa yang diukur} di bawah model multisumbu yang formal. CVX
membawa vaksin. Satu kelompok \emph{code system} HL7 lainnya membawa nilai administratif milik
FHIR sendiri---\emph{clinical status}, \emph{category}, \emph{encounter class}---yang
merupakan metadata struktural alih-alih terminologi klinis dan kelompok tersebut mencakup
kira-kira separuh dari seluruh \emph{code system} yang hadir.

Pembagian tersebut tampak langsung pada graf yang dimuat. Observasi menjangkau LOINC melalui
sisi \texttt{CODE} (242.217), tetapi menjangkau SNOMED~CT melalui sisi
\texttt{VALUE} (47.259): LOINC menamai pertanyaannya dan SNOMED~CT menamai
jawabannya bilamana jawaban itu sendiri berupa konsep berkode, bukan angka. Status merokok
merupakan contoh paling jelas---LOINC \texttt{72166-2} mengidentifikasi observasinya dan
nilainya berupa konsep SNOMED~CT.

Akibatnya, terdapat dua konsekuensi rancangan. Pertama, simpul
\texttt{:Code} harus memiliki kunci berupa (\texttt{system}, \texttt{code}), bukan berupa kode semata,
sebab untai kode saja dapat berbenturan antarkosakata. Kedua,
kosakata-kosakata itu merupakan \emph{namespace} terpisah yang dikelola organisasi berbeda
dengan siklus rilis yang mandiri dan tidak ada satu sisi pun yang menghubungkannya di dalam
terminologi sumber mana pun. Satu rekomendasi panduan klinis lazimnya menggunakan beberapa kosakata
sekaligus. Contohnya, rekomendasi mengenai terapi statin bagi pasien penyakit jantung koroner kronis menyebutkan satu bahan aktif
RxNorm dan satu gangguan SNOMED~CT; begitu ambang batas lipid diberlakukan, satu pengukuran LOINC juga
disebut. Tidak ada properti apa pun dalam data pasien yang dapat menghubungkan ketiganya. Karena itu, lapisan
konsep UMLS bukanlah sekadar kemudahan, melainkan keharusan struktural---lapisan itulah satu-satunya
artefak yang memberi ketiganya satu pengenal bersama dan tanpanya lapisan panduan klinis sama sekali
tidak dapat menjangkau lapisan pasien.

Konsekuensi ketiga menyangkut kontribusi penelitian ini secara khusus. LOINC memisahkan
pengukuran dari hasilnya dan menyisakan nilai numeriknya di dalam \texttt{valueQuantity}.
Seandainya fraksi ejeksi direkam sebagaimana SNOMED~CT merekamnya---sebagai temuan seperti
``disfungsi sistolik ventrikel kiri''---tidak akan ada angka yang dapat dibandingkan terhadap
sebuah ambang batas dan \emph{observation gate} pada Subbab~\ref{sec:dss-method} sama sekali tidak
dapat dituliskan. Perolehan \emph{observation gate} yang dilaporkan pada Bab~IV bertumpu pada
sifat LOINC ini secara spesifik, bukan pada pengodean terstruktur secara umum. Sifat itu
dimiliki pula oleh ekspor tabular Synthea, yang juga membawa kode LOINC beserta nilai
numeriknya (Subbab~\ref{sec:csv-gating}).

\subsubsection*{Apa yang hilang dari ekspor tabular terdatarkan}
\label{sec:flattening}
Bab~I melandaskan rancangan FHIR-native pada pembandingan terhadap ekspor tabular terdatarkan.
Karena Synthea mengeluarkan kohort yang sama dalam kedua format, pembandingan itu dapat
ditunjukkan, bukan sekadar dinyatakan sebagai klaim. Tabel~\ref{tab:flattening} memaparkannya untuk panel
tekanan darah.

\begin{table}[h]
\centering
\caption{Panel tekanan darah pada ekspor CSV dan pada graf FHIR, untuk kohort yang sama.}
\label{tab:flattening}
\begin{tabular}{lrl}
\hline
\textbf{LOINC} & \textbf{Baris CSV} & \textbf{Graf} \\
\hline
85354-9 panel tekanan darah & 0      & satu \texttt{:Observation} per pengukuran \\
8480-6 sistolik             & 52.583 & \texttt{:Component} dari Observation itu \\
8462-4 diastolik            & 52.583 & \texttt{:Component} dari Observation itu \\
\hline
\end{tabular}
\end{table}

Panel tekanan darah tersebut tidak dihasilkan pada ekspor CSV. Panel tersebut tidak memiliki
barisnya sendiri dan kedua anggotanya berubah menjadi baris yang saling lepas sehingga
memulihkan ``tekanan darah pasien ini 121/80'' menuntut perakitan ulang. Graf menyimpan
105.166 simpul \texttt{:Component}, yaitu $2 \times 52.583$---setiap anggota
terpelihara dan masing-masing masih melekat pada panel yang menaunginya.

Seberapa mahal perakitan ulang itu merupakan pertanyaan empiris dan Subbab~\ref{sec:csv-gating}
menjawabnya dengan membangun kembali \emph{observation gate} dari ekspor tersebut. Jawabannya perlu dinyatakan di sini agar pembandingan ini tidak terbaca lebih kuat daripada
yang sebenarnya. Ekspor itu menyertakan kolom \texttt{ENCOUNTER} yang terisi pada seluruh
baris observasi. Karena itu, kedua anggota satu pengukuran dapat dipertemukan kembali
melalui kunci rujukan yang tersurat, bukan melalui penanda waktu yang diandaikan
bersesuaian. Pada kohort
evaluasi, 17.177 pasang terakit kembali seluruhnya tanpa menyisakan satu baris pun yang tak
berpasangan. Yang hilang karenanya adalah identitas panelnya sebagai satu kesatuan, bukan
keterpulihan nilainya.

Perbedaan kedua lebih halus. Status
merokok muncul pada CSV sebagai 47.259 baris yang hanya membawa untai
tampilan (\emph{display string})---35.541 \texttt{Never smoked tobacco (finding)}, 11.587
\texttt{Ex-smoker (finding)}, dan 131 \texttt{Smokes tobacco daily
(finding)}---dengan \texttt{TYPE} tercatat sebagai \texttt{text}. \emph{Resource} FHIR membawa
untai yang sama di dalam \texttt{valueCodeableConcept}, tetapi bersamaan dengan \texttt{system}
\texttt{http://snomed.info/sct} dan \texttt{code} \texttt{266919005}. Ekspor itu mempertahankan
labelnya dan membuang kodenya; graf mempertahankan keduanya.

Tiga hal perlu dijelaskan di sini. Pertama, jembatan konsep mengekspor kode dari \emph{resource}
kondisi, obat, kategori \emph{care plan}, dan prosedur, tetapi tidak dari nilai observasi
sehingga \texttt{266919005} tidak pernah menjadi simpul \texttt{:Concept}. Kedua, \emph{observation gate} untuk merokok membaca properti \texttt{display} pada simpul
nilainya, bukan \texttt{code}-nya, sehingga status merokok dicocokkan sebagai untai di kedua
sisi. Ketiga, tidak ada satu pun pasien dalam kohort evaluasi yang merupakan perokok aktif.
Kelima rekomendasi yang digerbangi status merokok karenanya mencatat nol \emph{true
positive}, nol \emph{false positive}, dan nol \emph{false negative}, sebagaimana dilaporkan
pada Subbab~\ref{sec:ablation-class}.

Karena itu, status merokok merupakan contoh yang buruk untuk mendukung klaim bahwa kode tetap
terbawa saat diekspor dan di bahasan ini contoh tersebut memang tidak dipakai seperti demikian. Contoh tersebut tetap
disertakan karena \emph{jenis} kehilangan yang ditunjukkannya terwujud dan berlaku umum berupa tabel datar
tidak punya tempat untuk mencatat sistem kode asal sebuah nilai. Pada korpus ini kehilangan
tersebut kebetulan tidak mengakibatkan apa pun. Bagi jalur FHIR, dampak yang nyata justru datang dari panel tekanan darah di atas, yang nilai komponennya dibaca sebagai
angka dan memang mengubah hasil.

Kosakatanya sendiri tidak hilang. Setiap kolom
\texttt{CODE} bersandar pada kosakata yang ditetapkan oleh berkasnya, bukan yang dideklarasikan
pada barisnya: \texttt{conditions.csv} dan \texttt{procedures.csv} memakai SNOMED~CT,
\texttt{medications.csv} memakai RxNorm, \texttt{observations.csv} memakai LOINC,
\texttt{immunizations.csv} memakai CVX, dan kamus datanya menyatakan masing-masing hal itu
secara eksplisit \cite{synthea_csv_dict}. Kolom \texttt{SYSTEM} justru hadir tepat pada dua
berkas yang dokumentasinya menyebut kosakatanya \emph{typically} SNOMED~CT dan tidak hadir pada
tiga berkas yang kosakatanya dinyatakan kamus itu secara tegas. Karena itu, asimetrinya beralasan
dan pada tingkatan kolom utuh ekspor tersebut tidak kehilangan apa pun.

Yang tidak dapat diungkapkan oleh kosakata yang ditetapkan per berkas adalah kode yang
sistemnya berubah \emph{di dalam} satu kolom---dan itu persis terlihat pada kasus \texttt{VALUE} yang berada pada
observasi di atas. Satu kolom itu memuat angka bagi kolesterol LDL sekaligus konsep SNOMED~CT
bagi status merokok sehingga tidak ada satu sistem yang dapat secara konsisten dideklarasikan oleh
kolom pendamping dan kodenya pun dibuang demi untai tampilannya. FHIR melekatkan
\texttt{system} pada setiap \emph{coding} secara tersendiri sehingga tidak menghadapi kendala
serupa. Kehilangan itu merupakan konsekuensi dari bentuk tabularnya, bukan dari rancangan skema
yang terlewat.

Satu kehilangan lagi menyangkut ketelitian, bukan struktur. Ekspor tersebut menuliskan nilai
observasi ke satu angka di belakang koma, sedangkan \emph{resource} FHIR menyimpannya utuh.
Selisihnya terbatas pada $0{,}05$ dan berpotensi berarti hanya bila sebuah nilai kebetulan
berada tepat di tepi ambang panduan klinis. Subbab~\ref{sec:csv-gating} menghitung seberapa
sering hal itu terjadi pada kohort ini.

Yang benar-benar hilang secara struktural adalah \texttt{status}. Tidak satu pun tabel membawanya:
\emph{resource} FHIR mencatat \texttt{"status": "final"} dan CSV tidak memiliki tempat untuk
menaruhnya. Ketiadaan yang sama pada \texttt{medications.csv}, yang menyediakan \texttt{START},
\texttt{STOP}, dan \texttt{DISPENSES}, tetapi tanpa status, justru itulah yang harus dibahas
pada Subbab~\ref{sec:chronicity}.
Kedua total di bawah ini seluruhnya berasal dari proses umum: proses yang diperkaya penyakit
sama sekali tidak mencatat observasi tekanan darah maupun status merokok.

\subsubsection*{Mengapa observasi dibatasi dengan \emph{allow-list}}
Observasi merupakan penentu skala yang dominan dalam korpus FHIR: observasi mencakup sekitar
38\% dari seluruh \emph{resource} dan seorang pasien dalam kohort ini dapat
membawa beberapa ribu di antaranya. Memuat semuanya akan melipatgandakan graf beberapa kali
akibat adanya pengukuran rutin yang tidak menjadi prasyarat oleh satu pun rekomendasi
panduan klinis---misalnya urinalisis berulang, tanda vital dari kunjungan yang tidak berkaitan, dan
anggota panel yang tidak pernah dirujuk oleh logika yang dikodekan. \emph{Loader} karenanya
menerapkan \emph{allow-list} berisi sembilan kode LOINC yang persis mencakup nilai-nilai
yang menjadi tumpuan logika panduan klinis: fraksi ejeksi, indeks massa tubuh, kolesterol LDL dan
total, trigliserida, hemoglobin terglikasi, status merokok, serta panel tekanan darah. Tekanan
sistolik dan diastolik tidak memerlukan entri tersendiri karena keduanya dibawa sebagai
komponen di dalam panel dan ikut tertangkap bersamanya.

Keputusan ini berkaitan dengan \emph{skala} dan perlu dipisahkan dari mekanisme
\emph{observation gate} yang diuraikan pada Subbab~\ref{sec:dss-method}, yang merupakan
keputusan mengenai \emph{daya ungkap}. \emph{Allow-list} mengendalikan observasi mana yang ada di
dalam graf, sedangkan \emph{observation gate} mengendalikan bagaimana sebuah rekomendasi
yang bergantung pada nilai numerik diputuskan setelah observasi tersebut ada. Memperlebar
\emph{allow-list} akan memperbesar graf tanpa mengubah cara kerja penggerbangan dan mempersempitnya
akan membuat gerbang kekurangan masukan tanpa mengubah struktur logikanya.
\emph{Allow-list} ini diturunkan dari logika panduan klinis yang dikodekan dengan memperhatikan
aturan berupa sebuah kode akan dimasukkan karena ada rekomendasi yang membacanya. Karena itu, tidak ada rekomendasi yang dikeluarkan dari evaluasi akibat ketiadaan observasi yang tidak perlu---daftar
tersebut membatasi seberapa banyak korpus yang dimuat, bukan rekomendasi mana yang dapat
dinilai.

\subsection{Dua sudut pandang atas skema}
Graf yang dimuat akan diuraikan di sini menurut dua sudut pandang yang saling melengkapi karena skema
lengkap dan bagian skema yang benar-benar ditelusuri lapisan rekomendasi berbeda jauh
ukurannya. Skema \emph{lengkap} adalah himpunan seluruh tripel label-relasi-label yang ada
dalam basis data: 24 label simpul---tidak termasuk label struktural \texttt{:Resource} yang
melekat pada segala sesuatu kecuali \texttt{:Code}, akibat konsekuensi metode konstruksi---yang dihubungkan
oleh 95 tripel berbeda. Skema lengkap direproduksi pada Lampiran~\ref{app:schema}. Skema
\emph{terpakai} adalah subgraf yang ditelusuri kueri: sebelas label dan sebelas tipe relasi. Dua belas
label klinis tetap dimuat, tetapi tidak dibaca oleh satu pun kueri yang dilaporkan dalam
skripsi ini---\texttt{:AllergyIntolerance}, \texttt{:CarePlan}, \texttt{:CareTeam},
\texttt{:Device}, \texttt{:Encounter}, \texttt{:ImagingStudy}, \texttt{:Immunization},
\texttt{:Location}, \texttt{:MedicationAdministration}, \texttt{:Organization},
\texttt{:Practitioner}, dan \texttt{:PractitionerRole}---dan seluruhnya tetap tersedia untuk
penelitian lanjutan. Label ketiga belas yang tidak terkueri, \texttt{:Guideline}, tidak
terpakai karena alasan yang berbeda: label itu membawa provenans panduan klinis melalui
\texttt{(:Guideline)$-$[:CONTAINS]$\rightarrow$(:Recommendation)}, sedangkan kueri pendukung
keputusan memasuki lapisan panduan klinis langsung pada \texttt{:Recommendation} sehingga simpul
induknya tidak pernah diperlukan pada saat kueri dijalankan.

Subgraf terpakai tidak dipilih berdasarkan pertimbangan subjektif. Subgraf itu adalah himpunan
label dan tipe relasi yang muncul secara harfiah pada modul rekomendasi,
\emph{observation gate}, \emph{briefing}, dan kronisitas. Setiap sisi pada Gambar~\ref{fig:graph-model} telah diverifikasi keberadaannya di dalam basis data yang dimuat,
bukan disimpulkan dari kode sumber \emph{loader}.

Simpul \texttt{:Code} bersama merupakan poros bagi keseluruhan rancangan: setiap tipe \emph{resource}
klinis bertemu padanya dan simpul itu merupakan satu-satunya titik lekat bagi lapisan konsep
UMLS, yang digunakan untuk membawa masing-masing dari lapisan panduan klinis dan literatur. Inilah yang memungkinkan
seorang pasien terhubung ke sebuah rekomendasi dan bukti pendukungnya hanya dengan menggunakan
penelusuran, tanpa pencocokan teks apa pun pada saat kueri dijalankan. Perhatikan arah
jembatannya: sisi tersebut berjalan
\texttt{(:Concept)-[:SAME\_CONCEPT]->(:Code)}, dari konsep ke kode, sehingga penelusuran yang
berangkat dari pasien mencapai konsep dengan menyusurinya secara terbalik.

\begin{figure}[htbp]
\centering
% Scaled to the measure: the natural width is a few mm over \textwidth.
\resizebox{\textwidth}{!}{%
\begin{tikzpicture}[
    node distance = 3.5mm and 10mm,
    res/.style  = {rectangle, rounded corners=2pt, draw, text width=27mm,
                   align=center, minimum height=7mm, font=\scriptsize},
    hub/.style  = {rectangle, rounded corners=2pt, draw, very thick, text width=21mm,
                   align=center, minimum height=9mm, font=\scriptsize, fill=black!12},
    lay/.style  = {rectangle, rounded corners=2pt, draw, text width=23mm,
                   align=center, minimum height=9mm, font=\scriptsize},
    pat/.style  = {rectangle, rounded corners=2pt, draw, very thick, text width=17mm,
                   align=center, minimum height=9mm, font=\scriptsize, fill=black!6},
    lbl/.style  = {font=\tiny, inner sep=1pt},
    ar/.style   = {-latex}
]
% resource EHR
\node[res] (c)  {\texttt{:Condition}};
\node[res, below=of c] (p)  {\texttt{:Procedure}};
\node[res, below=of p] (m)  {\texttt{:MedicationRequest}};
\node[res, below=of m] (md) {\texttt{:Medication}};
\node[res, below=of md] (o) {\texttt{:Observation}};
\node[res, below=of o] (cm) {\texttt{:Component}};

% jangkar pasien
\node[pat, left=10mm of p, yshift=-8mm] (pt) {\texttt{:Patient}};

% jembatan terminologi dan lapisan atas
\node[hub, right=of p, yshift=-10mm] (code) {\texttt{:Code}\\(system, code)};
\node[lay, right=of code] (con) {\texttt{:Concept}\\(CUI UMLS)};
\node[lay, above right=6mm and 10mm of con] (rec) {\texttt{:Recommendation}\\(COR)};
\node[lay, below right=6mm and 10mm of con] (abs) {\texttt{:Abstract}\\(PMID)};

% SUBJECT: resource klinis berjangkar pada pasien
\draw[ar] (c) -- (pt);
\draw[ar] (p) -- (pt);
\draw[ar] (m) -- (pt);
\draw[ar] (o) -- (pt);
\node[lbl, above=0.5mm of pt] {\texttt{SUBJECT}};

% CODE / MEDICATION / COMPONENT
\draw[ar] (c)  -- (code);
\draw[ar] (p)  -- (code);
\draw[ar] (o)  -- (code);
\draw[ar] (cm) -- (code);
\draw[ar] (md) -- (code);
\draw[ar] (m)  -- (code);
\draw[ar] (o)  -- (cm);
\draw[ar] (m)  -- (md);

% jembatan berjalan Concept -> Code
\draw[ar] (con) -- node[lbl, above, sloped] {\texttt{SAME\_CONCEPT}} (code);
\draw[ar] (rec) -- (con);
\draw[ar] (abs) -- (con);
\draw[ar] (rec) -- node[lbl, right] {\texttt{CITES}} (abs);

\node[lbl, above=1mm of c]   {lapisan EHR};
\node[lbl, above=1mm of rec] {lapisan panduan klinis};
\node[lbl, below=1mm of abs] {lapisan literatur};
\end{tikzpicture}}
\caption{Skema \emph{terpakai}: sebelas label dan sebelas tipe relasi yang ditelusuri kueri pendukung
keputusan. \emph{Resource} klinis berjangkar pada \texttt{:Patient} melalui \texttt{SUBJECT}
dan bertemu pada simpul \texttt{:Code} bersama; \texttt{:Concept} menjembatani \emph{ke}
\texttt{:Code} melalui \texttt{SAME\_CONCEPT}, sekaligus membawa lapisan panduan dan
literatur. Obat diresolusi melalui dua lintasan---langsung ke \texttt{:Code} atau melalui
rujukan \texttt{:Medication}---dan kueri menyusuri keduanya. Sisi
\texttt{CO\_OCCURS\_IN\_LITERATURE}, yang dihilangkan di sini demi keterbacaan, turut
menghubungkan \texttt{:Code} dengan \texttt{:Code}. Skema lengkap berisi 95 tripel diberikan
pada Lampiran~\ref{app:schema}.}
\label{fig:graph-model}
\end{figure}

\section{Metode Rekomendasi}
\label{sec:dss-method}

Untuk seorang pasien, lapisan rekomendasi menjalankan empat operasi deterministik
terhadap graf:

\begin{enumerate}
    \item \textbf{Rekomendasi yang berlaku.} Sebuah rekomendasi dinyatakan berlaku jika pasien
    termasuk dalam kohort CCD (gerbang dasar) \emph{dan} pasien memenuhi seluruh target
    kondisi rekomendasi tersebut. Target kondisi adalah diagnosis lanjutan di luar gerbang
    kohort itu sendiri yang mungkin disyaratkan sebuah rekomendasi dan target itu hanya
    dihitung bila menjembatani ke kode yang benar-benar dipakai EHR. Penelusurannya bersifat
    gerbang-plus-AND, dengan OR untuk rekomendasi yang ditandai disjungtif. Subbab~\ref{sec:disc-class} memformalkan istilah
    tersebut.
    \item \textbf{\emph{Observation gate}.} Rekomendasi yang prasyaratnya bergantung pada
    observasi numerik (misalnya fraksi ejeksi, LDL, indeks massa tubuh, trigliserida, tekanan
    darah, dan status merokok) diputuskan dengan mengevaluasi predikat panduan klinis atas nilai
    terukur pasien yang dibaca dari graf, menggunakan ambang yang sama dengan standar rujukan
    (Subbab~\ref{sec:oracle}). Inilah mekanisme yang membuat observasi FHIR-native
    meningkatkan presisi.
    \item \textbf{Keamanan kontraindikasi.} Obat yang dikonsumsi pasien diperiksa terhadap
    rekomendasi yang mengontraindikasikannya dalam tiga varian. Varian unary menandai jika
    pasien mengonsumsi obat tersebut. Varian tergerbang-populasi menandai hanya jika pasien
    juga memenuhi kondisi populasi kontraindikasinya. Varian tergerbang-kronisitas, khusus
    untuk rekomendasi NSAID kronis, menandai hanya jika penggunaan obat memenuhi kriteria kronis. Gerbang kronisitas yang
    dipakai di dalam lapisan rekomendasi dihitung dari graf semata. Gerbang itu membaca
    \texttt{MedicationRequest.status}, yang oleh FHIR dicatat sebagai kategori
    (\texttt{active} atau \texttt{completed}) alih-alih sebagai tanggal, serta jumlah
    episode peresepan yang terpisah. Dengan demikian, gerbang ini hanya menjalankan bagian
    definisi kronisitas yang masih tersedia pada hasil ekspor FHIR. Karena tanggal berhenti tidak
    diekspor, kedua lengan dari definisi kronisitas yang dinyatakan sebagai durasi sama sekali tidak dapat
    dinilai dari graf; Subbab~\ref{sec:chronicity} memaparkan asimetrinya secara utuh.
    \item \textbf{Dukungan literatur dan \emph{briefing}.} Pilihan obat yang terikat dengan
    kondisi pasien pada lapisan literatur diperingkat menurut PMI dan seluruh fakta di atas
    dirakit menjadi muatan terstruktur yang kemudian dirumuskan oleh model bahasa besar
    menjadi ringkasan untuk klinisi. Model tersebut tidak pernah mengueri graf dan tidak
    pernah mengambil keputusan klinis; model hanya merumuskan ulang fakta hasil penelusuran
    yang membawa provenans.
\end{enumerate}

\subsection{Rancangan kueri dan prosedur \emph{observation gate}}
Subbab ini berupa sebuah rancangan kueri dan prosedur yang menghasilkan rekomendasi untuk menjawab RQ2.
Penerapan rekomendasi dievaluasi melalui satu penelusuran. Penelusuran diawali dari
pencarian id pasien, lalu kueri memeriksa apakah pasien berada dalam kohort CCD (konsep
gerbang dasar). Untuk setiap rekomendasi yang dapat diterapkan, pasien kemudian harus
memenuhi \emph{seluruh} target kondisi nonjangkar dari rekomendasi tersebut yang
menjembatani ke EHR. Syarat ini diimplementasikan sebagai logika AND, yaitu jumlah target
yang cocok harus sama dengan jumlah target yang dapat dijembatani. Untuk rekomendasi yang ditandai disjungtif, kueri memakai cabang OR sehingga
satu target yang cocok sudah memadai. Karena tipe pilihan
obat pada FHIR dapat dinyatakan baik secara langsung maupun melalui rujukan, kedua lintasan
obat (\texttt{:Code} langsung dan rujukan
\texttt{:Medication}$\to$\texttt{:Code}) akan ditambahkan untuk setiap obat yang muncul pada \texttt{:Medication} sehingga
obat yang dirujuk pasti terpetakan.

\subsubsection*{Bagaimana sebuah rekomendasi membedakan pasien}
\label{sec:disc-class}
Tidak setiap rekomendasi dapat mengerucutkan kohort. Begitu seorang pasien melewati gerbang CCD,
sebagian rekomendasi berlaku bagi pasien itu secara otomatis tanpa ada lagi yang perlu diuji.
Sebagian rekomendasi lainnya hanya berlaku bagi pasien yang memenuhi syarat tertentu.
Terdapat dua mekanisme persyaratan yang tidak saling bergantung untuk mengerucutkan kohort. 
Pertama, sebuah rekomendasi dapat membawa \emph{target kondisi} di luar konsep gerbang
yang menjembatani ke EHR sehingga penelusuran dapat membuang pasien yang tidak
memenuhi target kondisi. Kedua, rekomendasi dapat membawa \emph{prasyarat numerik} atas sebuah
observasi sehingga \emph{observation gate} dapat membuang pasien yang nilai observasi terukurnya (\emph{numerik})
tidak memenuhi ambang batas. Persilangan keduanya menghasilkan empat kelas, yang dirangkum pada
Tabel~\ref{tab:disc-class}.

\begin{table}[h]
\centering
\caption{Klasifikasi rekomendasi nonkontraindikasi menurut mekanisme yang tersedia untuk
mempersempit kohort.}
\label{tab:disc-class}
\begin{tabular}{lccl}
\hline
\textbf{Kelas} & \textbf{Target?} & \textbf{Numerik?} & \textbf{Rekomendasi} \\
\hline
Universal        & tidak & tidak & 38 \quad (NUTRI, PA, WM:1, APT:1---2 \ldots) \\
Tergerbang-nilai & tidak & ya    & \phantom{0}5 \quad (BB:1, BB:2, BPM:1, RAAI:1, RAAI:2) \\
Tercocok-konsep  & ya    & tidak & 15 \quad (APT:11---13, APT:4---5 \ldots) \\
Keduanya         & ya    & ya    & 10 \quad (LIPID:9, SGLT2:4, SGLT2:6, TP:2 \ldots) \\
\hline
\textbf{Total}   &       &       & \textbf{68} \\
\hline
\end{tabular}
\end{table}

Kelas \emph{universal} adalah yang terbesar. Rekomendasi nutrisi, aktivitas fisik, dan
pengelolaan berat badan termasuk di dalamnya karena panduan klinis hanya mensyaratkannya kepada setiap pasien
CCD; sistem yang benar adalah sistem yang memunculkan rekomendasi-rekomendasi tersebut bagi setiap anggota kohort. Rekomendasi
\emph{tercocok-konsep} membedakan melalui penelusuran---misalnya rekomendasi yang menyasar
pasien yang juga memerlukan antikoagulan oral hanya berlaku bila indikasi penggunaan obat tersebut terkode di EHR.
Rekomendasi \emph{tergerbang-nilai} membedakan pasien semata-mata melalui predikat.
Contohnya, \texttt{BB:1} dan \texttt{BB:2} bertumpu pada fraksi ejeksi, yang merupakan angka
terukur dan bukan diagnosis. Karena itu, penelusuran berbasis kondisi semata tidak dapat
memisahkan pasien yang dikenai rekomendasi tersebut. Kelas \emph{keduanya} bekerja
berdasarkan kedua mekanisme tersebut secara berurutan.

Akibatnya, terdapat dua konsekuensi. Pertama, dapat disalahpahami
bahwa rekomendasi universal tidak dapat menurunkan presisi, dengan alasan
bahwa rekomendasi itu berlaku bagi setiap anggota kohort menurut definisinya. Pernyataan tersebut
keliru dan Subbab~\ref{sec:ablation-class} mengukur seberapa keliru: kelas universal
menghasilkan sebagian besar \emph{false positive} sistem. Sebabnya, ``universal'' di sini
merupakan pernyataan tentang \emph{pengodean}, bukan tentang panduan klinisnya. Istilah universal berarti
rekomendasi tersebut tidak memiliki target kondisi yang dapat dikodekan dan dijembatani
ke kode EHR---bukan berarti teks panduan klinisnya tidak menetapkan prasyarat. Ketika pandua klinis memang
menetapkan prasyarat yang tidak pernah direpresentasikan pada lapisan pengetahuan, sistem
tidak memiliki cara untuk menahan rekomendasi tersebut dan standar rujukan yang membaca
logika panduan klinis dari rekam medis mentah menandai selisihnya sebagai \emph{false positive}.
Dengan demikian, batas utama presisi terletak pada \emph{seberapa banyak panduan yang berhasil
dikodekan}, bukan pada seberapa baik kueri mencocokkan pasien terhadap apa yang telah
dikodekan.

Kedua, ablasi \emph{observation gate} pada Subbab~\ref{sec:eval-method}
hanya dapat menyentuh 15 rekomendasi yang memang tergerbang: lima yang tergerbang-nilai
dan sepuluh pada kelas keduanya. Ke-53 rekomendasi sisanya dinilai dengan cara yang persis
sama, baik ketika penggerbangan dinyalakan maupun dimatikan. Jadi, perbedaan apa pun
yang ditunjukkan ablasi itu pasti berasal dari 15 rekomendasi tersebut dan tidak mungkin dari yang
lain. Hal itu masih sebuah prediksi, tetapi bukanlah asumsi dan karena itulah Bab~IV melaporkan keselarasan
kelas demi kelas---supaya dapat diperiksa kebenarannya. Hasilnya terbukti persis: pada kelas yang tidak tergerbang, adanya penggerbangan maupun tidak menghasilkan kesepakatan hingga ke cacah satuannya.

\subsubsection*{Mengapa \emph{observation gate} merupakan mekanisme terpisah}
Rekomendasi bergantung-observasi tidak dapat diselesaikan oleh penelusuran yang baru saja
diuraikan karena alasan daya ungkap, bukan karena skala. Kueri gerbang-plus-AND mencocokkan
\emph{konsep}: kueri itu menanyakan apakah rekam medis seorang pasien memuat kode yang
menjembatani ke konsep yang ditargetkan sebuah rekomendasi. Namun, prasyarat seperti ``LDL-C di
bawah 70~mg/dL'' atau ``fraksi ejeksi kurang dari atau sama dengan 40\%'' bukanlah sebuah
konsep---melainkan sebuah predikat berupa sebuah \emph{nilai}. Tidak ada pola untuk simpul dan
sisi yang dapat mengevaluasinya karena yang harus dibandingkan adalah sebuah angka yang
dibawa sebagai properti, terhadap sebuah ambang yang dinyatakan dalam teks panduan klinis dan sama
sekali tidak hadir di dalam graf.

Rekomendasi bergantung-observasi karenanya diputuskan secara terpisah: modul
\emph{observation gate} membaca fitur observasi gabungan milik pasien dari graf (nilai laboratorium
terburuk sepanjang riwayat, tekanan darah, dan status merokok terkini) dan mengevaluasi
predikat panduan klinis bagi setiap rekomendasi tersebut menggunakan ambang yang sama dengan standar
rujukan. Himpunan akhir rekomendasi yang dapat diberlakukan dihasilkan dari dua bagian. Pertama, rekomendasi
yang tidak tergerbang observasi mempertahankan apa pun yang telah diputuskan
penelusuran gerbang-plus-AND pada graf. Kedua, rekomendasi yang tergerbang justru diputuskan
oleh predikatnya dan keputusan itu \emph{menggantikan} keputusan graf, bukan diiriskan dengannya.
Bedanya terlihat setiap kali panduan klinis menyatakan prasyaratnya sebagai daftar pilihan. Ada satu rekomendasi yang berlaku jika pasien memiliki hipertensi, diabetes, penyakit ginjal
kronis, \emph{atau} fraksi ejeksi yang menurun. Mensyaratkan keempatnya sekaligus akan
membuang setiap pasien yang hanya memiliki salah satunya. Pasien-pasien tersebut tetap
terjaring hanya bila predikatnya yang memutuskan keberlakuan rekomendasi.

Gambar~\ref{fig:obs-gate} merangkum mekanisme tersebut sekaligus ablasi yang mengukurnya.

\begin{figure}[htbp]
\centering
\begin{tikzpicture}[
    box/.style  = {rectangle, rounded corners=2pt, draw, text width=30mm, align=center,
                   minimum height=12mm, font=\scriptsize},
    wide/.style = {rectangle, rounded corners=2pt, draw, text width=62mm, align=center,
                   minimum height=10mm, font=\scriptsize},
    ev/.style   = {wide, fill=black!8},
    ar/.style   = {-latex, thick}
]
\node[box] (pt) {Pasien yang lolos gerbang CCD};

\node[box, below=9mm of pt, xshift=-20mm] (trav)
    {Penelusuran gerbang-plus-AND\\[1pt]\emph{mencocokkan konsep pada graf}};
\node[box, below=9mm of pt, xshift=20mm] (feat)
    {Fitur observasi dari graf\\[1pt]\emph{LVEF, LDL, BMI, TG, tekanan darah, status merokok}};

\node[box, below=8mm of trav] (keep)
    {53 rekomendasi tak bergerbang\\[1pt]\emph{keputusan graf dipertahankan}};
\node[box, below=8mm of feat] (pred)
    {15 rekomendasi bergerbang\\[1pt]\emph{predikat menggantikan keputusan graf}};

\node[wide] (set) at ($(keep)!0.5!(pred) + (0,-20mm)$)
    {Himpunan rekomendasi yang dapat diberlakukan};
\node[ev, below=7mm of set] (abl)
    {Ablasi: lengan tanpa gerbang versus lengan dengan gerbang};

\draw[ar] (pt) -- (trav);
\draw[ar] (pt) -- (feat);
\draw[ar] (trav) -- (keep);
\draw[ar] (feat) -- (pred);
\draw[ar] (keep) -- (set);
\draw[ar] (pred) -- (set);
\draw[ar] (set) -- (abl);
\end{tikzpicture}
\caption{Mekanisme \emph{observation gate} beserta ablasinya. Penelusuran pada graf
mencocokkan \emph{konsep} sehingga prasyarat berupa ambang atas sebuah \emph{nilai} tidak
dapat dievaluasi olehnya. Rekomendasi bergantung-observasi karena itu diputuskan oleh
predikat tersendiri dan keputusan predikat tersebut menggantikan keputusan graf alih-alih
diiriskan dengannya---hal inilah yang membuat prasyarat disjungtif menambah pasien, bukan
menyaringnya keluar. Ablasi menjalankan alur yang sama sebanyak dua kali, dengan gerbang
dimatikan dan dinyalakan, sehingga hanya ke-15 rekomendasi bergerbang itulah yang dapat
berubah.}
\label{fig:obs-gate}
\end{figure}


\emph{Grounded briefing} dihasilkan dalam empat tahap. Tahap~(i) adalah penelusuran
deterministik untuk ringkasan pasien, rekomendasi yang berlaku, tanda keamanan, dan dukungan
literatur melalui kueri di atas. Tahap~(ii) merakit seluruhnya menjadi muatan fakta
terstruktur yang setiap butirnya membawa provenans, yaitu pengenal rekomendasi,
\emph{class of recommendation}, dan PMID. Tahap~(iii) menggunakan satu model bahasa untuk
merumuskan fakta tersebut; setiap rekomendasi hanya boleh dijustifikasi dari kondisi pemicu
hasil penelusurannya dan model dilarang menyatakan prasyarat yang tidak didukung data.
Tahap~(iv) mengeluarkan fakta dan ringkasan sebagai artefak terpisah agar
\emph{faithfulness}-nya pada tingkat sitasi dapat diperiksa secara mekanis.
Listing~\ref{lst:applicable-core} memuat inti kueri tersebut, sedangkan teks lengkapnya
disalin pada Lampiran~\ref{app:queries}.
\begin{lstlisting}[caption={Inti kueri rekomendasi yang dapat diberlakukan menggunakan gerbang CCD, yang kemudian dilakukan pencocokan
target menurut tanda konjungsi/disjungsi milik rekomendasi itu sendiri.},
label={lst:applicable-core}]
MATCH (p:Patient {uid:$pid})
WHERE EXISTS { (p)<-[:SUBJECT]-(:Condition)-[:CODE]->(:Code)
                 <-[:SAME_CONCEPT]-(:Concept {umlsCui:$gate}) }
MATCH (rec:Recommendation) WHERE NOT rec.cor STARTS WITH '3'
/* ... kumpulkan target non-jangkar tiap rekomendasi sebagai extraConds,
   dan yang benar-benar dipenuhi pasien sebagai matched ... */
WHERE size(extraConds) = 0
   OR (coalesce(rec.targetLogic,'AND') = 'OR'  AND size(matched) >= 1)
   OR (coalesce(rec.targetLogic,'AND') = 'AND' AND size(matched) = size(extraConds))
RETURN rec.recommendationId AS recId
\end{lstlisting}

Gambar~\ref{fig:briefing-pipeline} memperlihatkan susunan tersebut. Hal yang penting adalah
model bahasa berada \emph{di hilir} penelusuran dan tidak memiliki jalur kembali ke graf.
Model tidak memiliki cara untuk memperoleh rekomendasi di luar hasil penelusuran sehingga
rekomendasi apa pun di luar himpunan itu pasti karangan. Justru itulah yang membuat
pemeriksaan butir yang dapat disitasi dapat diputuskan. Susunan ini tidak menjamin bahwa
setiap kalimat yang dirumuskan model setia secara makna (Subbab~\ref{sec:faithfulness}).

\begin{figure}[htbp]
\centering
\begin{tikzpicture}[
   node distance = 8mm,
   box/.style = {rectangle, rounded corners=2pt, draw, text width=24mm, align=center,
                 minimum height=13mm, font=\scriptsize},
   ar/.style  = {-latex, thick}
]
\node[box] (g) {\emph{Knowledge graph}\\[1pt]\emph{penelusuran Cypher deterministik}};
\node[box, right=of g] (f) {Muatan fakta\\[1pt]\emph{recId, COR, PMID}};
\node[box, right=of f] (l) {Model bahasa\\[1pt]\emph{hanya merumuskan}};
\node[box, right=of l] (b) {Ringkasan\\[1pt]\emph{utama + rincian}};
\node[box, below=10mm of l, fill=black!8] (chk) {Pemeriksaan \emph{faithfulness}};

\draw[ar] (g) -- (f);
\draw[ar] (f) -- (l);
\draw[ar] (l) -- (b);
\draw[ar] (f) |- (chk);
\draw[ar] (b) |- (chk);
\end{tikzpicture}
\caption{Alur \emph{grounded briefing}. Penelusuran bersifat deterministik; model bahasa hanya
merumuskan fakta hasil penelusuran dan pemeriksaan \emph{faithfulness} membandingkan
keluarannya kembali terhadap muatan fakta yang sama.}
\label{fig:briefing-pipeline}
\end{figure}

\subsection{Kronisitas obat: status FHIR versus durasi CSV}
\label{sec:chronicity}
Satu rekomendasi---menghindari penggunaan NSAID kronis pada penyakit jantung koroner
kronis---bergantung bukan pada \emph{apakah} sebuah obat diresepkan, melainkan pada
\emph{berapa lama} obat itu dikonsumsi dan rekomendasi ini menyingkap asimetri
antara dua tampilan dari rekam medis sintetis yang sama. \texttt{MedicationRequest} pada FHIR membawa tanggal \texttt{authoredOn} dan \texttt{status}
penggunaan (\texttt{active} versus \texttt{completed}), tetapi \emph{tidak} membawa tanggal
berhenti. Sebaliknya, ekspor CSV Synthea yang berpasangan dengannya membawa tanggal
\texttt{START} dan \texttt{STOP} secara eksplisit, tetapi \emph{tidak} membawa informasi
status. Keduanya mengodekan fakta dasar yang sama melalui cara
yang berbeda. \texttt{STOP} yang kosong pada CSV (penggunaan masih berjalan) bersesuaian satu-satu
dengan \texttt{status} FHIR bernilai \texttt{active}. Tidak ada proyeksi yang merupakan
subhimpunan dari yang lain.

\subsubsection*{Mengoperasionalkan istilah ``kronis''}
Panduan tidak mendefinisikan istilah yang dapat dijadikan sebagai tumpuan kronisitas. Rekomendasi
\texttt{CCD-2023:APT:10} (COR~3-Harm, LOE~B-R) berbunyi lengkap: \emph{``In patients with CCD,
chronic nonsteroidal anti-inflammatory drugs should not be used because of increased
cardiovascular and bleeding complications.''} Rekomendasi itu tidak menyebutkan durasi, jumlah
penggunaan, maupun pola dosis. Setiap implementasi karenanya harus menggunakan definisi
operasionalnya sendiri dan definisi tersebut merupakan keputusan penelitian yang harus
diwujudkan untuk keperluan operasional, bukan fakta yang dapat disitasi. Hal ini sendiri merupakan temuan: sebuah
rekomendasi yang tidak ambigu bagi pembaca manusia ternyata tidak membuatnya langsung dapat dieksekusi
dan celahnya harus ditutup oleh pihak yang mengimplementasikan panduan klinis tersebut.

Tidak ada satu pun definisi terbitan yang mencakup istilah ini, tetapi ambang-ambang yang
diterbitkan pun tidak sembarangan. Ambang-ambang tersebut mengelompok. Pada jangka pendek, sebuah
tinjauan sistematis mengenai keamanan NSAID jangka pendek mengambil ``sepuluh hari atau
kurang'' dan mencatat bahwa \emph{Food and Drug Administration} Amerika Serikat menyarankan
konsultasi dengan klinisi bila penggunaan melampaui sepuluh hari \cite{aminoshariae2016} dan
sebuah tinjauan pakar mengenai osteoartritis menyebut periode penggunaan awal yang lazim adalah tujuh
hingga sepuluh hari dengan penilaian ulang sekitar tiga minggu
\cite{magni2021}. Di luar itu, ambang eksplisit muncul pada tiga titik. Sebuah
studi farmakovigilans memperlakukan paparan NSAID lebih dari 21~hari
berturut-turut sebagai penggunaan yang berpotensi tidak tepat atau berlebihan
\cite{ussai2015}. Sebuah studi populasi menghitung penggunaan NSAID dari
konsumsi yang dilaporkan sendiri selama sekurang-kurangnya 30 hari
\cite{hessami2025}. Sebuah studi \emph{cross-sectional} mengenai peresepan pada
lansia bahkan menyatakan secara langsung bahwa seorang pasien dapat dianggap sebagai pengguna NSAID
kronis bila mengonsumsi NSAID sekurang-kurangnya tiga bulan
\cite{abdu2020}. Studi lain menghindari jumlah hari penggunaan dan
mengklasifikasikan penggunaan kronis melalui \emph{medication possession ratio}
\cite{campbell2022} atau melalui persentase hari terpasok \cite{kikuchi2021}.

Definisi yang diadopsi dalam penelitian ini memiliki empat lengan. Dua ambangnya berada di
dalam kelompok terbitan yang baru diuraikan dan dipilih karena alasan tersebut; satu lengan sisanya
merupakan parameter yang dinyatakan tanpa preseden langsung:

\begin{enumerate}
    \item[(i)] periode penggunaan yang sedang berjalan (\texttt{active});
    \item[(ii)] paparan kumulatif melampaui 90~hari;
    \item[(iii)] tiga periode penggunaan terpisah atau lebih;
    \item[(iv)] satu periode penggunaan tunggal melampaui 30~hari.
\end{enumerate}

Lengan~(ii), paparan kumulatif di atas 90~hari, bersesuaian dengan batas tiga bulan yang
dinyatakan secara eksplisit sebagai penggunaan kronis oleh Abdu \textit{et al.}
\cite{abdu2020}. Ambang tunggal dijelaskan dalam literatur tersebut. Lengan~(iv), satu
periode penggunaan di atas 30~hari, berada di antara dua batas terbitan terdekat. Batas
pertama adalah 21 hari berturut-turut yang oleh Ussai \textit{et al.} sudah dianggap tidak
tepat \cite{ussai2015}. Ambang itu satu-satunya dalam kelompok ini yang dinyatakan
dalam hari \emph{berturut-turut} sehingga paling dekat konstruksinya dengan satu periode
penggunaan berkesinambungan. Batas kedua adalah tanda 30 hari yang digunakan Hessami
\textit{et al.} \cite{hessami2025}. Lengan~(iii), tiga periode penggunaan terpisah atau
lebih, tidak memiliki preseden langsung; metode berbasis rasio cakupan
\cite{campbell2022}\cite{kikuchi2021} menangkap intuisi yang sama mengenai
penggunaan berulang alih-alih sekali pakai, tetapi tidak dalam bentuk yang dapat dialihkan ke
jumlah episode peresepan. Lengan itu dinyatakan di sini sebagai parameter penelitian ini.

Dua kualifikasi perlu menyertai hal ini. Yang pertama, yang bersesuaian dengan
karya terbitan adalah \emph{besaran} angkanya, bukan cara angka itu diperoleh. Hessami
\textit{et al.} mengukur apa yang dilaporkan sendiri oleh pasien sepanjang waktu dan Abdu
\textit{et al.} mengklasifikasikan pasien alih-alih periode penggunaan, sedangkan lengan-lengan
di sini dihitung dari tanggal \texttt{START} dan \texttt{STOP} pada peresepan. Ambang terbitan
menjustifikasi angka yang dipilih dengan menjelaskan bahwa ambang tersebut tidak menunjukkan bahwa yang dihitung adalah hal
yang sama.

Yang kedua, aturan ini sengaja dibuat dengan hati-hati. Aturan ini bersifat \emph{spesifik, bukan
sensitif}: sistem akan cenderung melewatkan pengguna kronis daripada menandai orang yang sebenarnya
bukan pengguna kronis. Pasien pada hari kelima belas tidak tertangkap, padahal hari kelima
belas sudah melampaui titik ketika penggunaan jangka pendek tanpa pengawasan mulai tidak
dianjurkan \cite{aminoshariae2016} dan sudah mendekati ambang yang akan ditandai Ussai \textit{et al.} Pada rekomendasi COR~3 (Harm), melewatkan kasus adalah kegagalan yang
diterima di sini agar setiap tanda yang benar-benar dimunculkan sistem tidak dapat dibantah.
Penggerbangan yang diterapkan pada graf yang diuraikan di bawah juga mengalami pertukaran yang sama sehingga rekomendasi
ini diperlakukan dengan cara yang sama di sepanjang penelitian.

Yang terpenting, \emph{temuan} yang dilaporkan skripsi ini tidak bergantung pada angka-angka
tersebut. Ambang mana pun yang dipilih, lengan~(ii) dan~(iv) dinyatakan dalam durasi, sedangkan
lengan~(i) dan~(iii) tidak. Klaim bahwa proyeksi FHIR mendukung lengan~(i) dan~(iii), tetapi
tidak mendukung lengan~(ii) dan~(iv) menjelaskan apa yang berhasil dipertahankan pada proses
ekspor. Klaim itu tidak menjelaskan di mana ketidakmampuannya berdampak. Hanya angka \emph{recall} yang dilaporkan pada Bab~IV yang akan
bergeser bila ambangnya diubah.

Dari keempat lengan, hanya~(i) dan~(iii) yang dapat dihitung dari graf---lengan~(i) dari
\texttt{status} yang kategoris dan lengan~(iii) dari cacah simpul
\texttt{:MedicationRequest} NSAID. Lengan berbasis durasi, yaitu~(ii) dan~(iv), memerlukan
tanggal berhenti dan karenanya tidak terlihat oleh alat yang hanya membaca FHIR secara native. Standar
rujukan (Subbab~\ref{sec:oracle}) mengevaluasi keempat lengan dari ekspor CSV, yang tetap
merupakan data sintetis mentah dan bukan keluaran sistem rekomendasi. Sebaliknya,
penggerbangan pada graf yang tertanam pada lapisan rekomendasi dan \emph{briefing} hanya
mengevaluasi lengan~(i) dan~(iii). Kesenjangan di antara keduanya bukanlah jalan pintas
implementasi, melainkan akibat langsung dari apa yang dapat dipertahankan dari ekspor FHIR; dampaknya
terhadap \emph{recall} tanda keamanan dikuantifikasi pada Bab~IV.
Gambar~\ref{fig:chronicity} memaparkan keempat lengan dan menandai dua di antaranya yang
selamat dari proyeksi FHIR. Aturan ini diimplementasikan pada
\texttt{silver\_oracle.py} atas delapan zat aktif NSAID---ibuprofen, \emph{naproxen}, \emph{celecoxib},
\emph{diclofenac}, \emph{meloxicam}, \emph{indomethacin}, \emph{ketorolac}, dan \emph{piroxicam}---yang dicocokkan pada
deskripsi obat.

\begin{figure}[htbp]
\centering
\begin{tikzpicture}[
  node distance = 3mm,
  arm/.style = {rectangle, rounded corners=2pt, draw, text width=95mm, align=left,
                minimum height=8mm, font=\scriptsize, inner sep=4pt},
  hit/.style = {arm, fill=black!12, very thick},
  miss/.style= {arm, dashed}
]
\node[hit]  (a1) {(i)~~Penggunaan sedang berjalan \hfill dapat dihitung: \texttt{status = active}};
\node[miss, below=of a1] (a2)
    {(ii)~~Paparan kumulatif $>$ 90 hari \hfill perlu tanggal berhenti};
\node[hit,  below=of a2] (a3)
    {(iii)~~Tiga periode penggunaan terpisah atau lebih \hfill dapat dihitung: cacah \emph{request}};
\node[miss, below=of a3] (a4)
    {(iv)~~Satu periode penggunaan $>$ 30 hari \hfill perlu tanggal berhenti};

% legenda
\node[draw, very thick, fill=black!12, minimum width=5mm, minimum height=2.5mm,
      inner sep=0pt, anchor=west] (k1) at ([yshift=-7mm] a4.west) {};
\node[font=\scriptsize, right=1mm of k1] (k1t) {dapat dihitung dari graf FHIR};
\node[draw, dashed, minimum width=5mm, minimum height=2.5mm,
      inner sep=0pt, right=6mm of k1t] (k2) {};
\node[font=\scriptsize, right=1mm of k2] {memerlukan ekspor CSV};
\end{tikzpicture}
\caption{Empat lengan aturan NSAID kronis. Hanya lengan (i) dan (iii) yang dapat dihitung dari
graf FHIR, sedangkan lengan berbasis durasi memerlukan tanggal berhenti yang tidak dibawa ekspor FHIR.}
\label{fig:chronicity}
\end{figure}

\section{Metode Evaluasi}
\label{sec:eval-method}

Evaluasi memadukan penilaian intrinsik terhadap artefak graf dengan penilaian ekstrinsik dari
keluaran rekomendasi.

\subsection{Standar rujukan: \emph{silver standard} yang divalidasi terhadap \emph{gold}}
\label{sec:oracle}

Evaluasi menggunakan dua himpunan rujukan dan pembedaan di antara keduanya mengikuti
pembedaan yang lazim dalam literatur penambangan teks biomedis. \emph{Gold standard} adalah
rujukan yang dikurasi secara \emph{manual}, sedangkan \emph{silver standard} adalah rujukan yang
dihasilkan secara \emph{otomatis}. Kedua istilah tersebut diperkenalkan bersama oleh
Rebholz-Schuhmann \textit{et al.} \cite{rebholz2010,rebholz2011}, yang membandingkan
``korpus \emph{gold standard} beranotasi manual'' dengan korpus yang dihasilkan secara
otomatis, dengan alasan dari dilakukannya pembuatan silver standard berdasar bahwa penyiapan yang pertama
``memakan waktu dan mahal'' sehingga tidak terskala hingga beberapa ribu butir.

Alasan tersebut juga digunakan di sini. Melabeli setiap pasien dalam kohort yang dapat
dievaluasi gerbang secara manual akan menuntut seorang manusia menilai kecocokan dari 876
pasien terhadap 81 rekomendasi---berkisar tujuh puluh ribu keputusan
satuan---yang tidak mungkin dikerjakan oleh seorang peneliti tunggal. Evaluasi karenanya mengadopsi rancangan dua tingkat yang lazim pada \emph{phenotyping} rekam
medis elektronik. Pada rancangan itu, label \emph{silver} diturunkan secara otomatis dalam
skala besar, lalu divalidasi terhadap himpunan \emph{gold} kecil yang dikurasi manual
\cite{nogues2022}. Di sini, sebuah \emph{silver standard} melabeli seluruh kohort,
sebagaimana diuraikan pada sisa subbab ini. Sampel berlabel manual kemudian berperan sebagai
\emph{gold standard} untuk memeriksa \emph{silver standard} tersebut
(Subbab~\ref{sec:goldlabel}).

Satu kualifikasi mengenai peristilahan. Pada karya korpus yang memperkenalkannya, sebuah
\emph{silver standard} dibangun dengan \emph{mengharmonisasi keluaran beberapa sistem otomatis
yang saling bebas} sehingga kesepakatan antarsistem menggantikan pertimbangan manusia.
\emph{Silver standard} yang digunakan di sini bukan artefak konsensus semacam itu, melainkan
merupakan satu pengodean aturan yang deterministik terhadap panduan klinis. Istilah tersebut karenanya
digunakan dalam pengertiannya yang lebih luas dan kini lazim---yaitu standar rujukan yang
dihasilkan secara otomatis---sebagaimana diterapkan pada literatur \emph{phenotyping} klinis
\cite{nogues2022}, bukan dalam pengertian harmonisasi yang spesifik seperti pada korpus CALBC asli.
Istilah \emph{oracle} yang dipakai di sepanjang skripsi ini merujuk pada implementasi
\emph{silver standard} tersebut. \emph{Oracle} adalah komponen yang menghitung label bagi setiap
pasangan pasien--rekomendasi dengan menerapkan logika panduan klinis pada rekam medis mentah,
melalui jalur yang terpisah dari sistem yang diuji. Dengan demikian, \emph{oracle},
\emph{silver standard}, dan standar rujukan otomatis menunjuk komponen yang sama.

\subsubsection*{Konstruksi dan dua tingkatan ketakbergantungan}
\emph{Silver standard} mengodekan logika panduan klinis secara langsung dan melabeli setiap pasien
dengan membaca \emph{bundle} FHIR \emph{mentah}---tidak pernah membaca graf maupun keluaran
sistem rekomendasi. Ketakbergantungan itu perlu dinyatakan secara terpilah sebab tidak semua
rekomendasi memperolehnya pada tingkatan yang sama dan menyebutnya sebagai satu sifat tunggal
akan menyatakan lebih daripada yang sebenarnya berlaku.

Bagi 53 rekomendasi yang tidak tergerbang-observasi, kedua sisi berbeda baik pada logikanya
maupun pada jalur datanya. Sistem memutuskan melalui penelusuran graf, yaitu gerbang dan
konjungsi atas kode yang dijembatani konsep UMLS, sedangkan \emph{oracle} memutuskan melalui
pencocokan untai atas nama kondisi yang dibacanya dari \emph{bundle}. Keduanya kerap tidak
bersepakat dan justru ketidaksepakatan itulah yang membuat presisi dan \emph{recall} bermakna
di sana.

Bagi 15 rekomendasi yang tergerbang-observasi, ketakbergantungannya hanya berlaku pada jalur
datanya. Ambang numerik dan agregasi milik \emph{oracle}---misalnya nilai laboratorium terburuk
sepanjang riwayat dan tekanan darah terkini---digunakan apa adanya oleh \emph{observation
gate} sehingga predikat pada kedua sisi bukan sekadar serupa, melainkan merupakan fungsi yang
sama. Hal itu diperiksa, bukan diandaikan. Kedua predikat dievaluasi atas vektor ciri yang sama,
yaitu 876 vektor kohort yang sesungguhnya dan 200.000 vektor sintetis yang sengaja diarahkan
ke tiap ambang yang diujinya. Tidak ditemukan satu pun ketidaksepakatan
(Lampiran~\ref{app:modules}, \texttt{predicate\_audit.py}).

Konsekuensinya perlu dinyatakan dengan jelas. Bagi kelima belas rekomendasi tersebut,
kesepakatan antara sistem dan standar rujukan tidak dapat menjadi bukti bahwa ambangnya benar
secara klinis sebab kedua sisi memang menerapkan ambang yang sama. Yang dibuktikannya adalah
bahwa \emph{loader} membawa nilai pasien dari \emph{bundle} ke graf tanpa mengubahnya---sebuah
klaim mengenai jalur datanya, bukan mengenai aturannya. Pengujian yang tidak bergantung
terhadap aturannya sendiri menuntut pertimbangan klinisi yang berpraktik. Telaah klinisi
(Subbab~\ref{sec:clinician-method}) memberikan pengujian semacam itu, tetapi hanya bagi tiga
dari kelima belas rekomendasi tersebut. Selebihnya dicatat sebagai keterbatasan pada
Subbab~\ref{sec:limitations}.
Subbab~\ref{sec:shared-predicate} melaporkan seberapa besar bagian angka keselarasan yang
bertumpu pada kelima belas rekomendasi tersebut.

Satu
pengecualian yang terdokumentasi adalah rekomendasi NSAID kronis, yang lengan kronisitas
berbasis durasinya dibaca dari ekspor CSV Synthea---yang tetap merupakan data sintetis mentah,
bukan keluaran sistem---karena ekspor FHIR tidak memuat tanggal berhenti obat;
asimetri ini beserta konsekuensinya dirinci pada Subbab~\ref{sec:chronicity}.

Modul yang mengimplementasikannya bernama \texttt{silver\_oracle.py} dan menuliskan
\texttt{eval/silver\_oracle\_labels.csv}, sedangkan label yang dikurasi manual tetap berada pada
\texttt{eval/gold\_labels.csv}; dengan demikian kedua nama berkas itu membawa pembedaan yang
sama dengan yang diuraikan pada teks. Lampiran~\ref{app:modules} mencantumkan modul-modul tersebut beserta
keluaran masing-masing.

\subsection{Evaluasi keselarasan panduan klinis dan ablasi \emph{observation gate}}
Metrik utama dari evaluasi ini adalah keselarasan panduan klinis: himpunan rekomendasi dari hasil sistem rekomendasi
dibandingkan dengan label \emph{oracle} pada kohort yang dapat dievaluasi gerbang,
dilaporkan sebagai presisi, \emph{recall}, dan $F_1$. Eksperimen utamanya adalah sebuah
\emph{ablasi} yang menjalankan sistem dengan dan tanpa \emph{observation gate}; selisihnya
mengisolasi kontribusi \emph{observation gate} itu sendiri, bukan kontribusi wadah datanya
(Subbab~\ref{sec:csv-gating}). Setiap rekomendasi selanjutnya
diklasifikasikan menurut seberapa jauh ia membedakan pasien (apakah targetnya terkodekan dan
menjembatani ke EHR), yang menjelaskan sumber \emph{false positive} yang tersisa.

\subsection{Pelabelan \emph{gold} manual dan kesepakatan \emph{silver}-\emph{gold}}
\label{sec:goldlabel}
Untuk memvalidasi \emph{silver standard}, sampel berisi 30 pasien dilabeli secara
manual terhadap teks panduan klinis dan rekam medis lengkap pasien, menghasilkan \emph{gold
standard}. \emph{Silver standard} kemudian dibandingkan terhadap himpunan \emph{gold} tersebut
menggunakan Cohen's $\kappa$ \cite{cohen1960}, sebuah statistik kesepakatan terkoreksi-peluang (\emph{chance-corrected agreement})
yang juga telah digunakan untuk membandingkan keluaran model bahasa terhadap label rujukan
pada konteks klinis \cite{shah2024}. Ketidakpastian $\kappa$ ditaksir dengan \emph{bootstrap}
\cite{efron1994} yang mencuplik ulang \emph{pasien}, bukan sel. Sel-sel dari satu pasien tidak
saling bebas karena rekam medis yang sama menentukan seluruh labelnya dan pencuplikan ulang
per klaster mempertahankan ketergantungan tersebut \cite{field2007}. Karena $\kappa$ juga
dipengaruhi prevalensi label dan bias antarpenilai, indeks prevalensi, indeks bias, dan PABAK
dilaporkan bersamanya \cite{byrt1993}. Kesepakatan \emph{silver}-\emph{gold} yang tinggi itulah
yang mengizinkan penggunaan \emph{silver standard} atas seluruh kohort. Kesepakatan tersebut
menetapkan bahwa label yang dihasilkan otomatis mereproduksi pertimbangan manusia pada sampel
tempat keduanya tersedia. Dengan demikian, angka keselarasan yang dihitung atas seluruh
876 pasien bertumpu pada rujukan yang tervalidasi, bukan rujukan yang
sekadar diasumsikan.
Pelabelan dilakukan oleh penulis sendiri, dengan membaca teks panduan klinis terhadap rekam
medis lengkap tiap pasien. Empat keputusan cukup sering berulang sehingga ditetapkan di
muka. Indikasi antikoagulasi mencakup fibrilasi atrium, katup mekanis atau prostetik,
tromboemboli vena, serta resep warfarin yang sedang berjalan. Terapi antiplatelet ganda
mencakup aspirin bersama clopidogrel, \emph{ticagrelor}, atau \emph{prasugrel}. Sindrom koroner akut
mencakup STEMI, NSTEMI, dan sindrom prainfark. Intensitas statin mengikuti batasan yang ditetapkan panduan klinis. Atorvastatin~40---80~mg
dan \emph{rosuvastatin}~20~mg tergolong intensitas tinggi. Atorvastatin~10---20~mg,
\emph{rosuvastatin}~5---10~mg, \emph{simvastatin}~20---40~mg, \emph{pravastatin}~40---80~mg, serta
ezetimibe/\emph{simvastatin}~40~mg tergolong intensitas sedang. \emph{Simvastatin}~10~mg,
\emph{pravastatin}~20~mg, serta \emph{lovastatin}~10~mg tergolong intensitas rendah. Cacah rekomendasi yang dilabeli tidak seragam antarpasien, yaitu 71 hingga
73 rekomendasi, dengan keseluruhan 2.160 sel. Penyebabnya, tiga rekomendasi kontraindikasi
(\texttt{APT:7}, \texttt{APT:9}, dan \texttt{APT:10}) hanya dilabeli bagi pasien yang
memang mengonsumsi obat bersangkutan sebab bagi pasien lain rekomendasi itu tidak
memiliki subjek yang dapat dinilai. Karena
pelabelnya tunggal, keandalan antarpenilai bagi \emph{gold standard} itu sendiri tidak
dapat dilaporkan; hal ini ditinjau kembali di antara keterbatasan pada
Subbab~\ref{sec:limitations}.

\subsection{Evaluasi deteksi keamanan}
Deteksi keamanan kontraindikasi dievaluasi terhadap tanda keamanan yang dilabeli manual.
Varian unary dan varian tergerbang-populasi masing-masing dinilai presisi dan
\emph{recall}-nya. Kedua nilai itu memperlihatkan pertukaran antara penandaan berlebih pada
varian unary dan terlewatnya kontraindikasi yang secara struktural tidak dapat digerbangi
pada varian tergerbang-populasi. Rekomendasi NSAID kronis dinilai secara terpisah di bawah varian
tergerbang-kronisitasnya, yang berdampak pada perubahan nilai \emph{recall} akibat bergantung pada status obat
kategoris FHIR ketika lengan kronisitas berbasis durasi tidak tersedia
(Subbab~\ref{sec:chronicity}).

\subsection{Evaluasi \emph{faithfulness} ringkasan}
\emph{Grounded briefing} dievaluasi \emph{faithfulness}-nya melalui prosedur otomatis yang
membandingkan teks yang dihasilkan terhadap fakta hasil penelusuran. Prosedur itu menandai
setiap pengenal rekomendasi atau PMID yang dikarang, setiap rekomendasi berlaku atau tanda
keamanan yang terlewat, serta satu pola klaim berlebih, yaitu pernyataan bahwa pasien
memerlukan antikoagulan oral padahal tidak ada antikoagulan pada daftar obatnya. Sebuah
ringkasan dinyatakan lolos bila keempat hitungan itu nol. Prosedur ini tidak memeriksa apakah
setiap kalimat lainnya didukung fakta sebab pemeriksaan semacam itu memerlukan pembacaan
manusia.

\subsection{Telaah ketepatan rekomendasi oleh klinisi}
\label{sec:clinician-method}
Tingkat-tingkat evaluasi di atas menguji sistem terhadap aturan yang diturunkan dari panduan
klinis. Telaah klinisi menguji hal yang tidak dapat diuji aturan itu sendiri, yaitu apakah
rekomendasi yang dikeluarkan sistem sesuai bagi pasien yang bersangkutan menurut penilaian
klinisi. Telaah ini menilai keluaran sistem, bukan label \emph{gold standard}.

\subsubsection*{Lembar telaah}
Lembar telaah memuat 20 pasien, yaitu 20 pasien berskor tertinggi di antara 30 pasien sampel
\emph{gold standard} menurut aturan pemilihan sampel tersebut (Subbab~\ref{sec:goldlabel}).
Setiap pasien disajikan pada satu halaman yang berisi identitas singkat, kesimpulan
\emph{grounded briefing}, dan tabel rekomendasi. Kesimpulan itulah satu-satunya dasar
penilaian. Vinyet klinis dan daftar obat lengkap sengaja tidak ditampilkan sehingga penilai
menilai apa yang benar-benar disampaikan sistem, bukan rekam medis di baliknya.

Kesimpulan tersebut diterjemahkan ke dalam bahasa Indonesia. Terjemahannya diperiksa secara
mekanis: sitasi, target tekanan darah, ambang indeks massa tubuh, stadium, kelas rekomendasi,
dan nama obat harus terbawa. Penulis kemudian menyunting 12 dari 20 terjemahan. Suntingan itu
menghapus sitasi dan penanda kelas rekomendasi di dalam kurung, menghapus frasa ``yang sudah
tegak'', mengembalikan beberapa istilah ke bentuk Inggrisnya, dan merumuskan ulang beberapa
frasa. Penilai karenanya menilai teks yang telah disunting, bukan keluaran model apa adanya.

Tabel rekomendasi memuat seluruh tanda keamanan bagi pasien tersebut dan seluruh rekomendasi
berlaku yang dipicu kondisi pasien. Perawatan CCD baku yang berlaku sama bagi setiap pasien
tidak dimuat. Keseluruhannya berjumlah 327 baris dari 21 jenis rekomendasi dalam tiga
kelompok, yaitu cek/stop penggunaan obat (107 baris), tata laksana nonobat (77 baris), dan terapi
medikamentosa (143 baris). Kode dan kelas rekomendasi tidak ditampilkan. Keduanya disimpan
pada berkas kunci yang tidak dikirim kepada penilai. Setiap baris dinilai dengan satu dari tiga
pilihan, yaitu ``Tepat'', ``Belum tepat'', atau ``Tidak dapat dinilai''. Panduan penilaian
menyatakan pilihan ketiga sebagai jawaban yang sah sebab kesimpulan tidak selalu memuat
seluruh informasi yang diperlukan. Penilai diminta menuliskan saran pada baris yang tidak
dinilai tepat.

\subsubsection*{Penilai dan cakupan}
Tiga klinisi menilai secara terpisah dan diminta tidak berdiskusi sebelum penilaiannya
selesai.
Penilai~A dan penilai~C merupakan dokter spesialis jantung dan pembuluh darah (SpJP), dengan
penilai~C sebagai yang lebih senior. Penilai~B merupakan dokter umum yang sedang menempuh
program doktor. Tidak seorang pun menyelesaikan seluruh lembar. Penilai~A menilai pasien
1---15, penilai~B menilai pasien 1---10, dan penilai~C hanya menilai pasien~1. Di dalam
cakupan tersebut, penilai~A dan penilai~B masing-masing melewatkan satu baris. Jumlah baris
yang dinilai karenanya 247 bagi penilai~A, 160 bagi penilai~B, dan 16 bagi penilai~C.

\subsubsection*{Analisis}
Proporsi setiap pilihan dihitung per penilai. Interval kepercayaan 95\% bagi proporsi
``Tepat'' ditaksir dengan \emph{bootstrap} 10.000 replikasi yang mencuplik ulang pasien,
dengan alasan yang sama seperti pada $\kappa$ \emph{silver}-\emph{gold}
(Subbab~\ref{sec:goldlabel}). Kesepakatan penilai~A dan penilai~B diukur pada baris yang
dinilai keduanya dengan kesepakatan teramati, Cohen's $\kappa$ untuk tiga kategori
\cite{cohen1960}, dan PABAK \cite{byrt1993}, beserta interval \emph{bootstrap} per pasien.
Penilaian penilai~C hanya mencakup satu pasien sehingga dilaporkan secara deskriptif. Saran
dibaca dan dikelompokkan menurut rekomendasinya. Seluruh perhitungan dilakukan oleh
\texttt{baca\_hasil.py} dan \texttt{analisis\_hasil.py} (Lampiran~\ref{app:modules}).

\subsection{Evaluasi struktural graf}
\label{sec:structural}
Evaluasi sejauh ini bersifat \emph{ekstrinsik}: evaluasi itu menilai graf melalui mutu
keputusan yang diambil darinya. Graf juga dinilai secara \emph{intrinsik}, sebagai artefak
tersendiri, karena hasil rekomendasi tidak mengatakan apa pun mengenai apakah
representasi yang mendasarinya cukup heterogen, cukup luas, atau cukup dalam keterhubungannya
untuk dapat digunakan ulang.

Penilaian ini mengadopsi kerangka evaluasi ontologi yang ditetapkan oleh Gangemi \textit{et
al.} \cite{gangemi2006}. Kerangka itu membedakan tiga jenis ukuran. Ukuran \emph{struktural}
memperlakukan ontologi sebagai graf. Ukuran fungsional menyangkut kelayakan ontologi bagi
penggunaan yang dituju dan ukuran \emph{usability-profiling} menyangkut anotasinya. Tingkat
yang diuraikan di sini bersifat struktural dalam pengertian tersebut, sedangkan tingkat
keselarasan, keamanan, \emph{faithfulness}, dan telaah klinisi di atas bersifat fungsional. Metrik
struktural dari keluarga ini telah dievaluasi secara empiris atas repositori ontologi publik
dan tidak satu pun di antaranya menghasilkan pengelompokan yang tidak stabil: atas 119
ontologi OBO~Foundry dan 78 ontologi AgroPortal, klasifikasi yang ditimbulkan tiap metrik
bertahan terhadap \emph{bootstrap resampling} korpusnya dan dalam pengertian itulah
metrik-metrik tersebut disebut stabil \cite{franco2020}.
Argumen yang sama dengan yang diajukan skripsi ini---bahwa cacah skala semata tidak
mencirikan sebuah \emph{knowledge graph}---juga mendasari penggunaan himpunan metrik struktural yang
diusulkan untuk \emph{knowledge graph} lintas domain \cite{seo2022}.

Tiga dimensi dilaporkan. \textbf{\emph{Richness}} menangkap heterogenitas skema dan kosakata:
jumlah label simpul, tipe relasi, dan sistem kode serta jumlah lapisan semantik yang
disatukan. Dimensi ini mengikuti gagasan \emph{richness} yang diperkenalkan OntoQA
\cite{tartir2005}, yang mencirikan sebuah ontologi melalui keragaman relasi dan kelas yang
benar-benar terpakai alih-alih melalui ukuran mentahnya. OntoQA merumuskan gagasan tersebut
sebagai rasio---\emph{relationship richness}-nya adalah $P/(SC+P)$, yaitu porsi relasi nonpewarisan di antara seluruh sambungan antarkelas---yang tidak dapat diterapkan pada graf
ini karena tidak ada hierarki \emph{subsumption} sehingga tidak ada $SC$ sebagai
pembaginya. Skripsi ini karenanya melaporkan keragamannya itu sendiri dalam bentuk cacah. \textbf{\emph{Width}} menangkap
keluasan cakupan: kardinalitas simpul dan sisi, ukuran kohort, rerata jumlah kondisi per
pasien, dan proporsi panduan klinis yang terepresentasikan. \emph{Width} dalam pengertian ini
bersesuaian dengan dimensi \emph{completeness} pada mutu \emph{knowledge graph}, yang dalam hal ini
tersedia literatur metodologis yang cukup luas \cite{issa2021,xue2023}.
\textbf{\emph{Depth}} menangkap seberapa jauh sebuah kesimpulan dapat dilacak: panjang lintasan
inferensi kanonis pasien-rekomendasi-bukti serta sejauh mana simpul
\texttt{:Code} bersama berperan sebagai simpul pusat.

Meskipun demikian, terdapat dua penyimpangan dari literatur sumber yang harus dinyatakan karena istilah yang digunakannya sama, tetapi
besaran yang terukur tidak. Pertama, \emph{depth} dalam literatur evaluasi ontologi umumnya berarti
kedalaman hierarki \texttt{is-a}---berapa banyak tingkat yang dimiliki sebuah pohon
kelas. Graf yang dibangun di sini adalah graf properti yang tidak memiliki hierarki
\texttt{is-a} untuk diukur sehingga \emph{depth} didefinisikan ulang sebagai panjang lintasan
inferensi
yang melintasi lapisan-lapisan yang disatukan. Kedua, \emph{breadth} pada ontologi umumnya
berarti jumlah kelas bersaudara pada satu tingkat, sedangkan \emph{width} di sini berarti
kardinalitas dan cakupan kohort. Jadi, yang dipinjam adalah \emph{penggolongan} ukurannya, bukan
rumusnya: pemisahan evaluasi struktural dari evaluasi fungsional pada Gangemi
\cite{gangemi2006} dan prinsip OntoQA yang mencirikan sebuah graf melalui keragaman relasinya
alih-alih melalui ukurannya \cite{tartir2005}. Ketiga besaran yang dilaporkan pada Bab~IV
didefinisikan oleh skripsi ini sendiri dan tidak semestinya dibandingkan secara angka terhadap
nilai yang diterbitkan di tempat lain dengan nama yang sama.

Cohen's $\kappa$ dan metrik keselarasan dihitung langsung dari tabel kesepakatan
$2\times2$ oleh skrip evaluasinya sendiri, tanpa pustaka statistik, sehingga
mereproduksinya tidak menuntut kebergantungan apa pun di luar pustaka standar Python.

\section{Etika, Masalah, dan Keterbatasan Penelitian}
\label{sec:ethics}

\subsection{Pertimbangan etis}
Data pasien yang digunakan dalam penelitian ini seluruhnya bersifat sintetis. Synthea
membangkitkan individu yang masuk akal secara statistik, tetapi fiktif \cite{synthea} sehingga
tidak ada informasi kesehatan yang perlu dilindungi yang diproses, tidak diperlukan persetujuan pasien,
dan tidak ada risiko identifikasi ulang. Hal ini merupakan pilihan rancangan yang disengaja, bukan sekadar kemudahan. Pilihan itu
memungkinkan rekam medis lengkap, termasuk riwayat obat dan lintasan nilai laboratorium,
diterbitkan bersama kodenya untuk direproduksi. Penerbitan semacam itu mustahil dilakukan
dengan rekam medis nyata, yang umumnya masih berupa teks naratif tak terstruktur.

Satu komponen penelitian ini melibatkan manusia, yaitu telaah klinisi
(Subbab~\ref{sec:clinician-method}). Para klinisi berperan sebagai penilai ahli atas keluaran
sistem, bukan sebagai subjek penelitian. Penelitian ini mengumpulkan penilaian mereka atas
rekomendasi bagi pasien sintetis, bukan data mengenai diri mereka. Keikutsertaannya bersifat
sukarela dan identitas penilai diganti dengan kode pada naskah maupun data yang diterbitkan.
Karena itu, persetujuan etik untuk penelitian dengan subjek manusia tidak diperlukan. Tingkat
evaluasi lainnya bersifat otomatis: standar rujukan dibangkitkan secara terprogram, standar
emas dilabeli oleh penulis dari rekam medis sintetis, dan \emph{faithfulness} ringkasan
diperiksa melalui prosedur verifikasi otomatis. Konsekuensi rancangan tersebut terhadap
kekuatan pembuktian sistem dinyatakan secara eksplisit sebagai keterbatasan pada
Subbab~\ref{sec:limitations}.

Pertimbangan ketiga adalah peran yang dituju sistem ini. Sistem dirancang sebagai alat
\emph{rekomendasi} yang menampilkan rekomendasi panduan klinis dan tanda kontraindikasi untuk
dinilai oleh klinisi; sistem tidak meresepkan dan tidak divalidasi untuk penggunaan otonom.
\emph{Grounded briefing} secara eksplisit dibatasi untuk merumuskan fakta hasil penelusuran
alih-alih bernalar secara klinis dan setiap pernyataan yang dibuatnya membawa pengenal
rekomendasi serta \emph{class of recommendation} sehingga pembaca dapat mengauditnya terhadap
panduan klinis sumber.

\subsection{Masalah dan keterbatasan}
\label{sec:limitations}
Keterbatasan berikut melekat pada rancangan dan membatasi hasil penelitian ini:

\begin{enumerate}
    \item \textbf{Data sintetis.} Pasien Synthea dibangkitkan dari model perjalanan penyakit
    sehingga pola komorbiditas dan lintasan nilai laboratoriumnya konsisten secara internal,
    tetapi tidak diambil dari populasi nyata. Kinerja yang diukur di sini karenanya menetapkan
    bahwa alur pemrosesan benar secara \emph{internal}, bukan bahwa alur itu akan mencapai
    keselarasan yang sama pada rekam medis nyata, yang lebih berderau, lebih jarang terisi,
    dan memuat teks tak terstruktur.
    \item \textbf{Tidak ada data gejala jantung.} Rekam medis Synthea memuat sejumlah ukuran
    yang dilaporkan pasien, yaitu skor nyeri serta kuesioner depresi, kecemasan, penggunaan
    zat, dan kebutuhan sosial. Namun, rekam medis itu tidak memuat gambaran gejala jantung. Pada
    30 pasien sampel, daftar keluhan utama setiap pasien tidak pernah berubah di sepanjang
    8.067 catatan klinisnya. Hanya satu pasien yang mencantumkan gejala jantung, yaitu rasa
    tertekan di dada dan sesak napas, dan tidak ada gradasi angina, frekuensi angina, maupun
    kapasitas latihan. Kelas fungsional NYHA tercatat sekali pada lima pasien, tetapi berada di
    luar \emph{allow-list} observasi. Tidak satu pun dari 876 pasien kohort membawa diagnosis
    angina berkode. Akibatnya, enam rekomendasi yang bergantung pada angina (\texttt{CR:1:2},
    \texttt{MTRA:1}---\texttt{MTRA:4}, dan \texttt{MRA:1}) tidak pernah bernilai positif
    sehingga bagian panduan klinis ini sama sekali tidak teruji. Pada rekam medis nyata, gejala
    semacam itu umumnya tercatat sebagai teks bebas dan menangkapnya menuntut pengodean
    terstruktur yang tidak dicakup rancangan ini.
    \item \textbf{Kehilangan informasi pada proyeksi FHIR.} Sebagaimana diuraikan pada
    Subbab~\ref{sec:chronicity}, ekspor FHIR membawa status obat yang kategoris, tetapi tidak
    membawa tanggal berhenti sehingga logika panduan klinis yang bergantung pada durasi tidak dapat
    dievaluasi dari graf semata. Hal ini merupakan sifat ekspornya, bukan sifat implementasinya,
    dan hal itu dikuantifikasi alih-alih disembunyikan.
    \item \textbf{Predikat bersama pada rekomendasi tergerbang-observasi.} Bagi 15 dari 68
    rekomendasi, predikat sistem dan aturan standar rujukan merupakan fungsi yang sama
    (Subbab~\ref{sec:oracle}). Kesepakatan sempurna pada kedua kelas tergerbang karenanya
    menguji jalur datanya, bukan ketepatan aturannya. Bagian yang dinilai penilai yang
    benar-benar tidak bergantung mencakup 53 rekomendasi dan mencatat presisi $0{,}666$ serta
    $F_1$ $0{,}797$, berbanding $0{,}685$ dan $0{,}811$ pada angka utama
    (Subbab~\ref{sec:shared-predicate}). Seluruh perolehan ablasi terukur pada kelima belas
    rekomendasi berpredikat bersama tersebut sehingga yang diukurnya adalah kesetiaan
    implementasi, bukan keabsahan klinis. Telaah klinisi hanya menguji aturan itu pada tiga
    dari kelima belas rekomendasi, yaitu \texttt{WM:2}, \texttt{WM:3}, dan \texttt{SGLT2:4}.
    Seluruh 50 penilaian penilai~A dan penilai~B atas ketiganya ``Tepat''
    (Subbab~\ref{sec:clinician}).
    \item \textbf{Pelabel \emph{gold standard} tunggal.} Seluruh 2.160 sel \emph{gold
    standard} dilabeli oleh satu orang, yaitu penulis, yang juga mengodekan aturan
    \emph{silver standard}. Karena itu, keandalan antarpenilai bagi \emph{gold standard} tidak
    dapat dilaporkan. Kesepakatan yang tinggi juga sebagian dapat mencerminkan tafsir yang sama
    terhadap teks panduan klinis, bukan semata ketepatan aturannya. Interval kepercayaan pada
    Subbab~\ref{sec:kappa-ci} menjawab ketidakpastian akibat ukuran sampel, tetapi tidak menjawab
    keterbatasan ini sebab pencuplikan ulang tidak dapat menghadirkan penilai kedua.
    Pelabelan ulang oleh penilai kedua yang tidak bergantung pada penulis merupakan
    satu-satunya cara untuk menutupnya.
    \item \textbf{FHIR bukan syarat bagi \emph{observation gate}.} Perolehan presisi yang
    dilaporkan Subbab~\ref{sec:ablation} berasal dari penggerbangan atas nilai observasi
    numerik, bukan dari wadah FHIR secara khusus. Gerbang yang sama dapat dibangun kembali dari
    ekspor tabular Synthea dan menghasilkan keputusan yang sama persis pada seluruh 13.140
    keputusan yang diperbandingkan (Subbab~\ref{sec:csv-gating}). Pembedaan itu membatasi apa
    yang boleh diklaim rancangan FHIR-native ini: kelangsungan dan kemudahannya, bukan
    keterlaksanaannya. Perlu ditambahkan bahwa hasil tersebut berlaku bagi ekspor Synthea, yang
    menyertakan kunci rujukan \texttt{ENCOUNTER} pada setiap baris observasi; ekspor tabular
    yang tidak menyertakannya belum tentu bernasib sama.
    \item \textbf{Satu panduan klinis, satu penyakit.} Hanya panduan AHA/ACC CCD 2023 yang dikodekan
    dan hanya untuk penyakit jantung koroner kronis. Generalisasi ke panduan klinis atau kondisi lain
    sangat dimungkinkan mengingat rancangan jembatan konsepnya, tetapi tidak didemonstrasikan.
    \item \textbf{Gerbang kohort berkonsep tunggal.} Sebagaimana diuraikan pada
    Subbab~\ref{sec:gate}, gerbang CCD menerima satu konsep UMLS, sedangkan panduan klinis
    mendefinisikan penyakit jantung koroner kronis secara disjungtif. Enam puluh tujuh pasien
    ($7{,}1\%$ dari kohort yang didefinisikan secara klinis) yang bukti CCD-nya hanya berupa
    kode \emph{situation} riwayat infark atau riwayat revaskularisasi akibatnya tidak pernah
    dicocokkan dengan rekomendasi apa pun. Hal ini membatasi hasil yang dilaporkan pada
    subhimpunan yang konservatif dan, pada penerapan nyata, akan menyebabkan pasien-pasien
    tersebut sama sekali tidak menerima keluaran. Memperlebar gerbang menjadi disjungsi konsep
    merupakan perubahan kecil pada kueri dan menjadi butir pertama penelitian lanjutan.
    \item \textbf{Tidak ada validasi prospektif atau berbasis luaran.} Keselarasan diukur
    terhadap \emph{silver standard} yang dihasilkan otomatis dan \emph{gold standard} yang
    dikurasi manual, bukan terhadap luaran pasien. Tidak ada satu pun hasil di sini yang
    menetapkan bahwa mengikuti rekomendasi sistem memperbaiki mutu tindakan perawatan klinis.
    \item \textbf{Perumusan yang stokastik.} Kata-kata dalam ringkasan dihasilkan oleh model
    bahasa dan berubah antarproses sehingga \emph{faithfulness} diverifikasi per proses
    alih-alih dibuktikan sekali untuk selamanya. Penelusuran, sebaliknya, bersifat
    deterministik dan dapat direproduksi.
    \item \textbf{Cakupan pemeriksaan \emph{faithfulness}.} Pemeriksaan otomatis mencakup 30
    ringkasan dan menguji butir yang dapat disitasi, kelengkapan rekomendasi, serta satu pola
    klaim berlebih. \emph{Faithfulness} makna setiap kalimat tidak diuji. Karena itu, klaim
    yang didukung adalah bahwa ringkasan tidak mengarang rekomendasi maupun sitasi dan tidak
    menghilangkan rekomendasi, bukan bahwa setiap kalimatnya setia pada fakta. Uji klaim
    berlebih itu sendiri diperbaiki sesudah keluarannya terlihat
    (Subbab~\ref{sec:faithfulness}). Tiga puluh ringkasan juga masih terlalu sedikit untuk
    menaksir laju kesalahan secara presisi.
    \item \textbf{Telaah klinisi yang terbatas.} Standar rujukan, pemeriksaan deteksi keamanan,
    dan verifikasi \emph{faithfulness} menguji sistem terhadap \emph{aturan yang diturunkan dari
    panduan klinis}. Hanya telaah klinisi (Subbab~\ref{sec:clinician}) yang mengujinya
    terhadap pertimbangan klinisi dan cakupannya terbatas. Penilai~A, satu-satunya SpJP yang
    menilai lebih dari satu pasien, mencakup 15 dari 20 pasien. Penilai~B, seorang dokter umum,
    mencakup 10 pasien dan penilai~C hanya pasien~1. Kesepakatan penilai~A dan penilai~B
    tergolong cukup ($\kappa = 0{,}29$) sehingga proporsi rekomendasi yang dinilai tepat
    bergantung pada siapa penilainya. Penilai juga hanya membaca kesimpulan dan tabel
    rekomendasi, bukan rekam medis lengkap. Sebagian penilaian ``Belum tepat'' dan ``Tidak
    dapat dinilai'' menyebut informasi yang tidak ditampilkan, misalnya fraksi ejeksi, status
    gagal jantung, risiko perdarahan, dan daftar obat. Kesimpulan yang dibaca penilai pun
    merupakan terjemahan yang telah disunting penulis. Hasil telaah ini karenanya menunjuk letak
    kelemahan sistem, tetapi belum menaksir selisih antara keselarasan dengan panduan klinis
    dan ketepatan klinis secara umum. Telaah ini juga menilai ketepatan rekomendasi, bukan mutu
    ringkasan. \emph{Reader study} dengan model Van Veen \textit{et al.}~\cite{vanveen2023},
    yang di dalamnya klinisi menilai ringkasan pada aspek kelengkapan, ketepatan, dan
    keringkasan dengan kesepakatan antarpenilai dikuantifikasi melalui Cohen's $\kappa$
    \cite{cohen1960}, masih diserahkan kepada penelitian selanjutnya.
\end{enumerate}
